\section{Virtio Over Messages}\label{sec:Virtio Transport Options / Virtio Over Messages}

\newcommand{\conceptref}[1]{\hyperref[sec:Virtio Transport Options / Virtio Over Messages / Basic Concepts / #1]{#1}}
\newcommand{\msgref}[1]{\hyperref[sec:Virtio Transport Options / Virtio Over Messages / Transport Messages / VIRTIO_MSG_#1]{VIRTIO_MSG_#1}}
\newcommand{\busref}[1]{\hyperref[sec:Virtio Transport Options / Virtio Over Messages / Bus Messages / BUS_MSG_#1]{BUS_MSG_#1}}
\newcommand{\msgdef}[1]{\subsubsection{VIRTIO_MSG_#1}\label{sec:Virtio Transport Options / Virtio Over Messages / Transport Messages / VIRTIO_MSG_#1}}
\newcommand{\busdef}[1]{\subsubsection{BUS_MSG_#1}\label{sec:Virtio Transport Options / Virtio Over Messages / Bus Messages / BUS_MSG_#1}}

This section defines \textbf{virtio-msg}, a transport mechanism that encapsulates
virtio operations as discrete message exchanges rather than relying on PCI or
memory-mapped I/O regions. It separates bus-level functionality (e.g., device
enumeration, hotplug events) from device-specific operations (e.g., feature
negotiation, virtqueue setup), ensuring that a single, generic transport layer
can be reused across multiple bus implementations.

virtio-msg addresses several key objectives:

\begin{itemize}
  \item \textbf{Support multiple bus implementations:}
  Systems can rely on various communication methods such as hypercalls, local
  IPC, network channels, or device trees for enumerating devices. virtio-msg
  defines a common transport interface suitable for any of these mechanisms.

  \item \textbf{Reduce per-bus complexity:}
  Buses can implement a fully message-based workflow (including optional
  enumeration via \busref{GET_DEVICES} and hotplug via \busref{EVENT_DEVICE})
  or they can discover and manage devices through
  alternative means such as platform firmware data. In either case, they
  forward transport messages to and from each device.

  \item \textbf{Preserve virtio semantics:}
  The transport leverages standard virtio concepts (features, configuration
  space, virtqueues), so existing virtio drivers and device implementations can
  integrate smoothly once a device is discovered.
\end{itemize}

\subsection{Basic Concepts}
\label{sec:Virtio Transport Options / Virtio Over Messages / Basic Concepts}

The virtio-msg transport relies on a set of foundational concepts to ensure
reusability across different bus implementations and flexibility in device
capabilities. This section defines those concepts and clarifies how they apply,
regardless of whether the bus leverages message-based enumeration or platform
data for device discovery.

\subsubsection{High-Level Architecture}

virtio-msg operates around two layers:

\begin{enumerate}
  \item \textbf{Bus Layer}: A \emph{bus instance} exposes zero or more virtio
        devices. It discovers available devices through mechanisms such as:
    \begin{itemize}
      \item Optional message-based queries (\busref{GET_DEVICES}),
      \item External data sources (e.g., device tree, ACPI tables, hypervisor
        firmware calls),
      \item Dynamic hotplug notifications (optionally via \busref{EVENT_DEVICE}).
    \end{itemize}
    Once a device number is identified, regardless of discovery method, the bus
    layer exposes that device number and availability state to the driver side
    of the transport.

  \item \textbf{Transport Layer}: After the bus knows about a device, the
        virtio-msg transport handles all device-specific operations:
    \begin{itemize}
      \item Retrieving and setting features (\msgref{GET_DEVICE_FEATURES},
        \msgref{SET_DRIVER_FEATURES}),
      \item Accessing the device configuration space (\msgref{GET_CONFIG},
        \msgref{SET_CONFIG}),
      \item Setting up virtqueues (\msgref{SET_VQUEUE}, \msgref{RESET_VQUEUE}),
      \item Managing status and runtime notifications (\msgref{SET_DEVICE_STATUS},
        \msgref{EVENT_USED}, etc.).
    \end{itemize}
    These transport messages remain the same across different bus instances,
    allowing a single virtio-msg driver component to function in multiple
    environments.
\end{enumerate}

\subsubsection{Relationship Between Bus and Transport}
\label{sec:Virtio Transport Options / Virtio Over Messages / Basic Concepts / Relationship between bus and transport}

This subsubsection explains the division of responsibilities: the bus layer is
the mandatory carrier that moves messages between driver and device endpoints
(e.g., over IPC, shared memory with signaling or hardware messaging), while
the virtio-msg transport defines the semantics of those messages.

virtio-msg groups messages into two categories:

\begin{description}
  \item[\textbf{Bus Messages}:]
    Intended for global bus operations such as enumerating device numbers
    (\busref{GET_DEVICES}), managing device hotplug events (\busref{EVENT_DEVICE})
    or checking peer responsiveness (\busref{PING}).
    These standard bus messages are provided for convenience. Support for
    each message is \emph{optional}. A bus binding selects which messages it
    uses and may define alternative mechanisms for these functions.

  \item[\textbf{Transport Messages}:]
    Used for device-specific operations, such as:
    \begin{itemize}
      \item Retrieving or setting features (\msgref{GET_DEVICE_FEATURES},
        \msgref{SET_DRIVER_FEATURES}),
      \item Accessing device configuration (\msgref{GET_CONFIG},
        \msgref{SET_CONFIG}),
      \item Managing virtqueues (\msgref{SET_VQUEUE}, \msgref{RESET_VQUEUE}),
      \item Handling device status and notifications (\msgref{SET_DEVICE_STATUS},
        \msgref{EVENT_USED}, etc.).
    \end{itemize}
\end{description}

This separation lets a bus remain minimal if it obtains device information from
firmware tables, while still supporting fully message-based enumeration and
hotplug when desired.

\busnormative{\paragraph}{Bus Message Implementation}{Virtio Transport Options / Virtio Over Messages / Basic Concepts / Relationship between bus and transport / Bus Messages}
\begin{itemize}
  \item A bus implementation MAY support any subset of \busref{GET_DEVICES},
        \busref{PING}, and \busref{EVENT_DEVICE}, including none of them.
  \item If a bus implementation supports one of these bus messages, it MUST
        conform to the normative definition for that message.
\end{itemize}

\busnormative{\paragraph}{Transport Message Forwarding}{Virtio Transport Options / Virtio Over Messages / Basic Concepts / Relationship between bus and transport / Transport Message Forwarding}
\begin{itemize}
  \item A bus implementation MUST validate the device number in each transport
        message before forwarding it.
  \item A bus implementation MUST relay each transport message to the device
        number identified in the message header, regardless of how it
        discovered or enumerated that device.
  \item A bus implementation SHOULD treat transport messages as opaque apart
        from enforcing generic transport limits, such as the advertised maximum
        message size, and SHOULD NOT modify the transport payload.
  \item If a bus implementation cannot validate or route a transport request
        that expects a response, it MUST surface a transport-visible failure
        for that request.
  \item If a bus implementation cannot validate or route a one-way transport
        event, it MAY drop the event or handle it according to
        message-specific bus requirements.
\end{itemize}

\subsubsection{Known Bus Implementation Examples}
\label{sec:Virtio Transport Options / Virtio Over Messages / Basic Concepts / Known Bus Implementation Examples}

This subsubsection lists informative references to documented bus bindings for
virtio-msg.

\begin{itemize}
  \item \textbf{Virtio Message Bus over FF-A}: specified by Arm in
        \hyperref[intro:FFA-BUS]{[FFA-BUS]}.
\end{itemize}

\subsubsection{System Topology}

A virtio-msg system contains the following elements:

\begin{itemize}
  \item \textbf{Bus Instances and Devices}: Each bus instance advertises its
        capabilities (e.g., transport revision, maximum message size) and
        discovers devices via message-based queries or external data sources.
        Every discovered device has a unique \emph{device number}.
  \item \textbf{Driver}: Communicates with the bus to learn about
        available devices. Once a device is recognized, the driver uses the
        common transport messages to perform feature negotiation, configuration,
        and virtqueue setup.
  \item \textbf{Device}: Implements virtio device functionality and
        responds to the transport messages. The bus forwards these messages to
        the correct device instance based on device number.
\end{itemize}

\subsubsection{Transport Revisions and Maximum Message Size}
\label{sec:Virtio Transport Options / Virtio Over Messages / Basic Concepts / Revisions}

Each \textbf{virtio-msg bus instance} advertises the following to the transport
layer:
\begin{itemize}
  \item \textbf{Transport revision}: the protocol version supported by the bus
        instance (independent of the overall Virtio specification version).
  \item \textbf{Maximum message size}: the largest message size (in bytes) that
        can be carried per request or response, including the common header.
  \item \textbf{Transport feature bits}: revision-specific optional features
        implemented by the bus instance.
\end{itemize}

Transport feature bits are determined per bus instance. When both sides expose
transport feature masks, the bus instance presents only the common subset as
the effective transport feature mask for that bus instance.

The mechanism for obtaining these parameters is implementation-defined and can
vary between bus instances. Common approaches include:
\begin{itemize}
  \item Reading firmware or device-tree data that describes each bus instance,
  \item Performing a message exchange during bus setup to retrieve the values,
  \item Relying on per-bus configuration structures or driver-specific defaults.
\end{itemize}

\busnormative{\paragraph}{Advertising Transport Parameters}{Virtio Transport Options / Virtio Over Messages / Basic Concepts / Revisions / Advertising Transport Parameters}
\begin{itemize}
  \item Each bus instance MUST make its transport revision, maximum message
        size, and transport feature bits available to the virtio-msg transport
        before any transport messages are exchanged for that device.
  \item A bus instance MUST expose transport feature bits independently of the
        device feature blocks exchanged by \msgref{GET_DEVICE_FEATURES} and
        \msgref{SET_DRIVER_FEATURES}.
  \item A bus instance MUST advertise only the transport feature bits supported
        by both sides of that bus instance.
  \item A bus instance MUST apply the same limits to both driver-originated and
        device-originated transport messages.
  \item For an active device, a bus instance MUST treat the advertised
        transport revision, maximum message size, and transport feature bits as
        stable.
  \item If these values change, a bus instance MUST notify the transport layer
        and MUST require transport-layer revalidation before accepting
        additional transport messages for the affected device.
\end{itemize}

\drivernormative{\paragraph}{Respecting Bus Limits}{Virtio Transport Options / Virtio Over Messages / Basic Concepts / Revisions / Driver Limits}
\begin{itemize}
  \item A driver MUST NOT send a transport message whose total size exceeds the
        maximum message size advertised for the target bus instance.
  \item A driver MUST NOT rely on transport features or messages that require a
        higher transport revision than the bus instance reports.
\end{itemize}

\devicenormative{\paragraph}{Respecting Bus Limits}{Virtio Transport Options / Virtio Over Messages / Basic Concepts / Revisions / Device Limits}
\begin{itemize}
  \item A device MUST ensure its responses and device-originated messages do
        not exceed the maximum message size advertised by the bus instance that
        relays them.
  \item A device MUST NOT require transport features or messages beyond the
        transport revision reported by the bus instance, and MUST respond with
        an error or ignore requests for unsupported features.
\end{itemize}

\paragraph{virtio-msg transport revisions}

The following table lists the currently defined virtio-msg transport revisions
and their protocol-defined maximum message sizes. Transport feature bits are
defined separately in
\ref{sec:Virtio Transport Options / Virtio Over Messages / Basic Concepts / Transport Feature Bits}.

\begin{tabular}{ |l|l|l| }
\hline
\field{revision} & \field{maximum size} & remarks \\
\hline \hline
1        & 52-65535    & Virtio Message Revision 1 \\
\hline
\end{tabular}

The table reflects the protocol-defined maximum size; the recommended
advertised maximum remains 264 bytes for interoperability.

Note that a change in the virtio standard does not necessarily
correspond to a change in the virtio-msg transport revision.

The maximum message size is specified from the transport-layer point of view
and includes the 8-byte common header plus payload. Any extra encapsulation
imposed by the underlying bus (for example, a framing header) does not count
against this limit. Today the largest practical inline configuration payload is
244 bytes: messages such as GET_CONFIG and SET_CONFIG use 12 bytes of
message-specific fields, leaving up to 244 bytes for device configuration data
in one transfer when the recommended 264-byte maximum message size is used.
Larger configuration regions can be accessed through multiple exchanges without
requiring a larger per-message limit.

\busnormative{\paragraph}{Message Size Bounds}{Virtio Transport Options / Virtio Over Messages / Basic Concepts / Revisions / Message Size}
\begin{itemize}
  \item A bus implementation MUST advertise a maximum message size of at least
        52 bytes.
  \item A bus implementation SHOULD NOT advertise a maximum message size that
        exceeds 264 bytes (256-byte payload plus the common header).
\end{itemize}

\paragraph{Versioning and Forward Compatibility}

A higher transport revision or additional transport feature bits extend the
protocol with new messages or capabilities. Interoperability across revisions
relies on retaining the mandatory messages and semantics of prior revisions,
as specified by the driver, device, and bus requirements below.

\drivernormative{\paragraph}{Revision Compatibility}{Virtio Transport Options / Virtio Over Messages / Basic Concepts / Versioning and Forward Compatibility / Driver}
\begin{itemize}
  \item A driver that negotiates transport revision $N$ MUST implement all
        mandatory driver behavior defined for revisions $1$ through $N$.
  \item If a driver receives a device-originated message or feature indication
        that requires an unsupported revision or transport feature, it MUST
        ignore the message (or treat the request as failed) and MUST NOT act on
        partially understood data.
\end{itemize}

\devicenormative{\paragraph}{Revision Compatibility}{Virtio Transport Options / Virtio Over Messages / Basic Concepts / Versioning and Forward Compatibility / Device}
\begin{itemize}
  \item A device that advertises transport revision $N$ MUST implement all
        mandatory device behavior defined for revisions $1$ through $N$.
  \item If a device receives a driver request that relies on an unsupported
        revision or transport feature, it MUST reject the request using the
        message-specific error mechanism (if any) or silently ignore it.
\end{itemize}

\busnormative{\paragraph}{Revision Compatibility}{Virtio Transport Options / Virtio Over Messages / Basic Concepts / Versioning and Forward Compatibility / Bus}
\begin{itemize}
  \item A bus instance that advertises transport revision $N$ MUST satisfy every
        bus requirement defined for revisions $1$ through $N$.
  \item If a bus instance cannot forward a message because it requires an
        unsupported revision or transport feature, it MUST surface a transport
        error or drop the message without forwarding it.
\end{itemize}

\subsubsection{Transport Feature Bits}
\label{sec:Virtio Transport Options / Virtio Over Messages / Basic Concepts / Transport Feature Bits}

Transport feature bits are bus-instance parameters. They are advertised as part
of the transport parameters described in
\ref{sec:Virtio Transport Options / Virtio Over Messages / Basic Concepts / Revisions}
and are not carried in the device feature blocks exchanged by
\msgref{GET_DEVICE_FEATURES} and \msgref{SET_DRIVER_FEATURES}.

The following transport feature bits are defined:

\begin{tabular}{|l|l|p{7cm}|}
\hline
Bit & Minimum revision & Meaning \\
\hline
0 & 1 & Select strict configuration semantics profile. \\
\hline
\end{tabular}

Bit 0 is named \texttt{VIRTIO\_MSG\_F\_STRICT\_CONFIG\_GENERATION}.
When this feature is not advertised for the bus instance, the baseline profile
applies. See
\ref{sec:Virtio Transport Options / Virtio Over Messages / Basic Concepts / Configuration Semantics Profiles}.

\subsubsection{Device Numbers and Enumeration}
\label{sec:Virtio Transport Options / Virtio Over Messages / Basic Concepts /
Device Numbers}

Each virtio-msg bus instance contains zero or more \emph{devices}, identified
by a 16-bit \textbf{device number}. The device number is a bus-assigned
identifier and is distinct from the virtio device ID (device type) returned by
\msgref{GET_DEVICE_INFO}. Buses discover these device numbers through
mechanisms such as:
\begin{itemize}
  \item \textbf{Message-Based Enumeration}: Using \busref{GET_DEVICES} to query
        which numbers exist (optional).
  \item \textbf{Platform Data}: A device tree, ACPI tables, or hypervisor calls
        might inform the bus of available device numbers and their properties.
\end{itemize}

The 16-bit device-number space provides roughly 65k device slots per bus
instance. Deployments that require more devices can instantiate multiple bus
instances and distribute devices across those instances.

Once a bus confirms that a device number is valid, regardless of the discovery
method, it exposes that number to the driver side of the transport. The driver
then issues \msgref{GET_DEVICE_INFO} as the first transport message for that
device number before other transport operations.

\busnormative{\paragraph}{Device Number Assignment}{Virtio Transport Options / Virtio Over Messages / Basic Concepts / Device Numbers / Assignment}
\begin{itemize}
  \item A bus implementation MUST assign a unique device number to every
        device on a given bus instance.
  \item A bus implementation MUST prevent requests, responses, or events
        associated with a removed device from affecting a subsequent device
        assigned the same device number.
  \item A bus implementation MUST provide the driver with sufficient
        information, either via \busref{GET_DEVICES}, \busref{EVENT_DEVICE},
        or equivalent platform data, to discover each valid device number.
\end{itemize}

\subsubsection{Configuration Semantics Profiles}
\label{sec:Virtio Transport Options / Virtio Over Messages / Basic Concepts / Configuration Semantics Profiles}

virtio-msg defines two configuration semantics profiles selected by the
transport feature bits advertised for the bus instance:
\begin{itemize}
  \item \textbf{Baseline profile}: applies when
        \texttt{VIRTIO\_MSG\_F\_STRICT\_CONFIG\_GENERATION} is not advertised for
        the bus instance.
  \item \textbf{Strict profile}: applies when
        \texttt{VIRTIO\_MSG\_F\_STRICT\_CONFIG\_GENERATION} is advertised for the
        bus instance.
\end{itemize}

The profiles apply to \msgref{GET_CONFIG}, \msgref{SET_CONFIG}, and
\msgref{EVENT_CONFIG}.

\drivernormative{\paragraph}{Configuration Semantics Profiles}{Virtio Transport Options / Virtio Over Messages / Basic Concepts / Configuration Semantics Profiles / Driver}
\begin{itemize}
  \item A driver MUST determine the active profile before issuing
        \msgref{SET_CONFIG}.
  \item A driver MUST determine the active profile from the transport feature
        bits advertised for the target bus instance, not from
        \msgref{GET_DEVICE_FEATURES}/\msgref{SET_DRIVER_FEATURES}.
  \item In any profile, a driver SHOULD use generation values from
        \msgref{GET_CONFIG} and \msgref{EVENT_CONFIG} to detect concurrent
        configuration changes and SHOULD re-read required configuration until
        it observes a stable generation value.
  \item In baseline profile, a driver MUST set \field{generation} to 0 in
        \msgref{SET_CONFIG} requests.
  \item In baseline profile, upon \msgref{EVENT_CONFIG}, a driver SHOULD refresh
        required configuration via \msgref{GET_CONFIG} when event data is
        omitted or when \field{offset} and \field{length} are zero.
  \item In strict profile, a driver MUST track generation values from
        \msgref{GET_CONFIG}/\msgref{EVENT_CONFIG} and MUST include its most
        recent generation value in \msgref{SET_CONFIG}.
  \item In strict profile, if a \msgref{SET_CONFIG} request is rejected due to a
        mismatched generation count, a driver SHOULD issue \msgref{GET_CONFIG}
        to obtain the latest configuration data and generation count before
        retrying.
  \item In any profile, a driver MUST NOT assume generation count is monotonic
        or preserved across device resets.
\end{itemize}

\devicenormative{\paragraph}{Configuration Semantics Profiles}{Virtio Transport Options / Virtio Over Messages / Basic Concepts / Configuration Semantics Profiles / Device}
\begin{itemize}
  \item A device MUST apply the configuration semantics profile selected by the
        transport feature bits advertised for the bus instance that relays its
        messages.
  \item In any profile, a device MUST provide its current generation count in
        \msgref{GET_CONFIG} responses, \msgref{SET_CONFIG} responses, and
        \msgref{EVENT_CONFIG} messages.
  \item In any profile, a device MUST change \field{generation} whenever there
        is a possibility that two accesses to the device configuration space
        can observe different versions of that space, as described in
        \ref{sec:Basic Facilities of a Virtio Device / Device Configuration Space}.
  \item In baseline profile, a device MUST ignore request \field{generation} in
        \msgref{SET_CONFIG}. For requests that are otherwise valid, the device
        MUST apply the write regardless of request \field{generation} and MUST
        NOT reject \msgref{SET_CONFIG} due solely to generation mismatch.
  \item A device MUST NOT send \msgref{EVENT_CONFIG} solely because a
        \msgref{SET_CONFIG} request carried a mismatched \field{generation}
        value.
  \item In any profile, a device MAY send a generic
        \msgref{EVENT_CONFIG} notification with \field{offset} and
        \field{length} set to 0 and no \field{data} when the affected
        configuration change cannot be represented as a bounded range
        (for example, when only an equivalent out-of-band notification is
        available without affected-range details).
  \item A device transport layer implementation MUST originate
        \msgref{EVENT_CONFIG} whenever it changes configuration data in a way
        that is visible to the driver.
  \item A device transport layer implementation MUST originate
        \msgref{EVENT_CONFIG} for status only when a device-originated
        asynchronous status change requires asynchronous reporting to the
        driver.
  \item A device MUST NOT send \msgref{EVENT_CONFIG} solely because the driver
        wrote status via \msgref{SET_DEVICE_STATUS}; such writes MAY be
        followed by \msgref{EVENT_CONFIG} only if a distinct
        device-originated asynchronous status change occurs afterward.
  \item If \msgref{EVENT_CONFIG} reports only a status change, a device MUST
        set \field{offset} and \field{length} to 0.
  \item In any profile, a device MAY reset generation count, including when it
        resets internal state.
  \item In strict profile, if a \msgref{SET_CONFIG} request includes a
        \field{generation} count that does not match the device's current
        count, the device MUST reject the request.
  \item If updated configuration data cannot fit in a single
        \msgref{EVENT_CONFIG}, a device SHOULD send an \msgref{EVENT_CONFIG}
        without configuration data and MUST make updated data available via
        \msgref{GET_CONFIG}. In strict profile, the device SHOULD set
        \field{offset}/\field{length} to the affected range.
\end{itemize}

This subsection defines baseline and strict semantics together so that
generation and mismatch behavior remains explicit and consistent across message
definitions.

\subsubsection{Device Feature Blocks}
\label{sec:Virtio Transport Options / Virtio Over Messages / Basic Concepts / Device Feature Blocks}

Device feature bits are organized into 32-bit blocks accessed by
\msgref{GET_DEVICE_FEATURES} and \msgref{SET_DRIVER_FEATURES}. Each block
represents up to 32 feature bits:

\begin{itemize}
  \item \textbf{Block Index}: The starting block (e.g., block 0 for
        features 0-31, block 1 for features 32-63, etc.).
  \item \textbf{Number of Blocks}: How many blocks the driver wishes to retrieve
        or modify in a single message.
  \item \textbf{Feature Data}: The 32-bit values representing either the
        device-offered feature bits or the driver-selected subset for the
        selected blocks.
\end{itemize}

Virtio-msg device feature blocks may include core virtio feature bits. They are
distinct from the transport feature bits defined in
\ref{sec:Virtio Transport Options / Virtio Over Messages / Basic Concepts / Transport Feature Bits}.
This transport does not support VIRTIO\_F\_NOTIF\_CONFIG\_DATA; device feature
negotiation requirements are defined in
\ref{sec:Virtio Transport Options / Virtio Over Messages / Device Initialization / Device Features}.

\devicenormative{\paragraph}{Device Feature Blocks}{Virtio Transport Options / Virtio Over Messages / Basic Concepts / Device Feature Blocks / Device}
\begin{itemize}
  \item When \msgref{GET_DEVICE_FEATURES} covers blocks that extend beyond the
        features a device implements, the device MUST return zero for the
        feature data in those positions.
  \item After processing \msgref{SET_DRIVER_FEATURES}, a device MUST update
        only the driver-selected feature blocks addressed by that request and
        MUST leave all other driver-selected feature blocks unchanged.
  \item A device MUST report the device-offered feature set in
        \msgref{GET_DEVICE_FEATURES} responses and MUST NOT reflect the
        driver-selected subset in those responses.
\end{itemize}

\subsubsection{Error Signaling}
\label{sec:Virtio Transport Options / Virtio Over Messages / Basic Concepts / Error Signaling}

Transport errors can arise from malformed messages, routing failures inside a
bus implementation, or faults while processing a valid request in the device
transport layer.
Implementations should handle such faults locally where possible. For
transport request/response messages that use the existing wire format, the
completion contract is deterministic:
\begin{itemize}
  \item For a valid, supported request accepted for processing, completion is
        exactly one protocol response.
  \item If routing, policy, timeout, or completion constraints prevent that
        response, completion is a transport-visible failure surfaced by the
        bus implementation.
  \item One-way events do not define a completion outcome. Unless a
        message-specific rule states otherwise, they MAY be dropped without a
        protocol error response or transport-visible failure.
  \item Malformed or unsupported transport messages are discarded without
        transport error responses; for request/response messages, the
        originating side observes completion according to the bounded
        transport-visible failure policy.
\end{itemize}

Bus bindings may define failure signaling exchanged only between bus
implementations and completion of underlying transfers. For a transport
request, this signaling can cause a local failure completion without waiting
for a timeout. Such signaling does not add transport messages or change the
request-completion contract.

\busnormative{\paragraph}{Error Handling}{Virtio Transport Options / Virtio Over Messages / Basic Concepts / Error Signaling / Bus}
\begin{itemize}
  \item A bus implementation MUST apply a bounded completion policy for each
        valid, supported transport request that expects a response.
  \item A bus implementation MUST complete each such request with exactly one
        of the following outcomes:
    \begin{itemize}
      \item one protocol response, or
      \item a transport-visible failure when routing, policy, timeout, or
            completion constraints prevent response delivery.
    \end{itemize}
  \item A bus implementation MUST NOT wait indefinitely for response delivery.
  \item A bus implementation MUST treat malformed headers or unsupported bus
        \field{msg_id} values in bus messages as invalid and MUST discard them.
        A bus binding MAY define bus-internal failure reporting or underlying
        transfer completion for such messages. A bus implementation MAY log
        the condition for diagnostics.
  \item A bus implementation MUST consume bus-specific failure signaling
        within the bus and MUST NOT deliver it to the transport as an inbound
        transport message.
  \item A bus implementation MUST NOT validate transport \field{msg_id} values;
        unsupported transport \field{msg_id} handling is performed by the
        transport layer.
  \item A bus implementation MUST NOT generate error responses to event
        (one-way) messages. This does not prohibit bus-internal completion of
        an underlying transfer.
  \item Unless a message-specific rule states otherwise, a bus implementation
        MAY discard an event without surfacing a completion outcome.
\end{itemize}

\drivernormative{\paragraph}{Error Handling}{Virtio Transport Options / Virtio Over Messages / Basic Concepts / Error Signaling / Driver}
\begin{itemize}
  \item A driver transport layer implementation receiving a malformed or
        unsupported transport message MUST discard it without producing a
        transport protocol response.
  \item A driver MUST apply a bounded timeout and retry policy when waiting for
        transport-request completion, including reset-completion checks.
  \item If that bounded policy is exhausted, a driver MUST treat the request as
        failed and MUST perform recovery using the reset and status semantics
        defined in
        \ref{sec:Virtio Transport Options / Virtio Over Messages / Device Operation}.
\end{itemize}

\devicenormative{\paragraph}{Error Handling}{Virtio Transport Options / Virtio Over Messages / Basic Concepts / Error Signaling / Device}
\begin{itemize}
  \item A device transport layer implementation receiving a malformed or
        unsupported transport message MUST discard it without producing a
        transport protocol response.
  \item Recovery actions taken in response to an error (such as retries,
        selective resets, or device removal) MUST follow the normative reset and
        status semantics defined in
        \ref{sec:Virtio Transport Options / Virtio Over Messages / Device Operation}.
\end{itemize}

This specification does not define a dedicated transport error-reporting
message. Transport-visible failure is surfaced as a local bus completion
outcome for a request.

Some bus implementations may reject requests that violate their local policy
or that are not valid for the current device state. In such cases, the bus may
report a local failure completion to the originating transport.

\subsubsection{Endianness}
\label{sec:Virtio Transport Options / Virtio Over Messages / Basic Concepts / Endianness}

Unless otherwise stated, all numeric fields defined for virtio-msg messages use
little-endian encoding.

\drivernormative{\paragraph}{Endianness}{Virtio Transport Options / Virtio Over Messages / Basic Concepts / Endianness / Driver}
\begin{itemize}
  \item A driver MUST encode and decode the common header and payload fields
        defined by this transport using little-endian byte order.
\end{itemize}

\devicenormative{\paragraph}{Endianness}{Virtio Transport Options / Virtio Over Messages / Basic Concepts / Endianness / Device}
\begin{itemize}
  \item A device MUST emit and consume the common header and payload fields
        defined by this transport using little-endian byte order.
\end{itemize}

\subsubsection{Memory Addresses}
\label{sec:Virtio Transport Options / Virtio Over Messages / Basic Concepts / Memory Addresses}

The address fields in \msgref{GET_VQUEUE}, \msgref{SET_VQUEUE}, and
\msgref{GET_SHM} use the address space and encoding defined by the bus binding.
These can be physical addresses or bus-specific representations of memory
locations. Memory alignment requirements apply to the identified memory, not
necessarily to the encoded address value.

\busnormative{\paragraph}{Memory Addresses}{Virtio Transport Options / Virtio Over Messages / Basic Concepts / Memory Addresses / Bus}
\begin{itemize}
  \item A bus implementation MUST define the address space and encoding used by
        the address fields in \msgref{GET_VQUEUE}, \msgref{SET_VQUEUE}, and
        \msgref{GET_SHM}, including any translation required to access the
        identified memory.
\end{itemize}

\subsubsection{Common Message Format}
\label{sec:Virtio Transport Options / Virtio Over Messages / Basic Concepts / Common Message Format}

All virtio-msg exchanges, whether \emph{bus messages} or \emph{transport
messages}, begin with an 8-byte header followed by an optional payload. The
fields below describe the wire format for that header.

The header layout is:
\begin{lstlisting}
struct virtio_msg_header {
    u8       type;      /* request/response + bus/transport */
    u8       msg_id;    /* message id */
    le16     dev_num;   /* target device number; bus messages MUST use 0 */
    le16     token;     /* bus-managed correlation identifier */
    le16     msg_size;  /* total size: header (8) + payload */
    u8       payload[];
};
\end{lstlisting}

Field semantics:
\begin{itemize}
  \item \field{type}:
    \begin{itemize}
  \item Bit[0]: 0=request, 1=response. Messages whose \field{msg_id} identifies
        an event MUST use Bit[0]=0 and MUST NOT be sent as responses.
  \item Bit[1]: 0=transport message, 1=bus message. This bit selects whether
        \field{msg_id} is interpreted in the transport-message namespace or the
        bus-message namespace.
  \item Bits[2..7]: reserved for future use.
    \end{itemize}
  \item \field{msg_id}: Message ID identifying the message definition. Ranges
    are defined in
    \ref{sec:Virtio Transport Options / Virtio Over Messages / Transport Messages}
    and
    \ref{sec:Virtio Transport Options / Virtio Over Messages / Bus Messages}.
    \begin{itemize}
  \item Bits[0..5]: message number within the selected range.
  \item Bit[6]: 0=request/response message, 1=event.
  \item Bit[7]: 0=standardized message ID, 1=implementation-defined message ID.
    \end{itemize}
  \item \field{dev_num}: For transport messages, the device number that should
    receive the message. Device number 0 is not reserved for transport
    messages. Bus messages are identified by \field{type} and MUST carry
    \field{dev_num}=0.
  \item \field{token}: Correlation identifier owned and managed by the bus.
    Drivers and devices treat this field as opaque. A bus implementation MAY
    transparently insert or overwrite \field{token} values on the request leg,
    response leg, or both to correlate concurrent requests and responses or to
    compensate for transport-specific reordering. A device copies the request
    \field{token} into the matching response before any bus-side rewrite.
  \item \field{msg_size}: Total size in bytes of the complete message (header
    plus payload).
  \item \field{payload}: Operation-specific data. If a bus introduces extra
    padding bytes, those bytes are not part of the payload semantics.
\end{itemize}

\drivernormative{\paragraph}{Common Header}{Virtio Transport Options / Virtio Over Messages / Basic Concepts / Common Message Format / Driver}
\begin{itemize}
  \item A driver MUST set bits 2..7 of \field{type} to zero and MUST treat them
        as reserved when parsing received headers.
  \item A driver MUST ensure \field{msg_size} reflects the total message length
        (header plus payload) and MUST NOT exceed the maximum message size
        advertised by the bus instance.
  \item When sending a transport message, a driver MUST set \field{dev_num} to
        the intended device number.
  \item If a driver introduces padding bytes that become part of the transport
        payload, it MUST set those bytes to zero.
  \item A driver MUST treat the \field{token} field as opaque and MUST NOT rely
        on its value when processing received messages.
\end{itemize}

\devicenormative{\paragraph}{Common Header}{Virtio Transport Options / Virtio Over Messages / Basic Concepts / Common Message Format / Device}
\begin{itemize}
  \item A device MUST set bits 2..7 of \field{type} to zero in transmitted
        messages and MUST ignore those bits on receive.
  \item A device MUST ensure \field{msg_size} reflects the total message length
        (header plus payload) and does not exceed the bus's advertised maximum.
  \item When sending a transport message, a device MUST set \field{dev_num} to
        its own device number.
  \item A device MUST ignore padding bytes that are documented as bus-specific
        and MUST zero any such bytes it introduces into the transport payload.
  \item A device MUST treat the \field{token} field as opaque and MUST NOT rely
        on its value when processing received messages.
  \item For each request/response message pair, a device MUST copy the
        \field{token} value from the received request into the matching
        response.
\end{itemize}

\busnormative{\paragraph}{Common Header}{Virtio Transport Options / Virtio Over Messages / Basic Concepts / Common Message Format / Bus}
\begin{itemize}
  \item A bus implementation MUST deliver bus messages with \field{dev_num}=0
        and MUST NOT alter \field{dev_num} for transport messages beyond the
        routing needed to reach the addressed device.
  \item A bus implementation MUST set bits 2..7 of \field{type} to zero when
        generating bus messages and MUST ignore those bits when forwarding
        transport messages.
  \item If the bus adds framing or padding bytes around the common header or
        payload, it MUST set those bytes to zero before delivering the message
        to the opposite side and MUST present the same zero padding when the
        opposite side reads a message.
  \item A bus implementation owns the \field{token} field; it MAY insert or
        overwrite values for correlation purposes and MUST ensure that any such
        use is transparent to drivers and devices.
  \item A bus implementation MAY rewrite \field{token} values on the request
        path, response path, or both, provided this remains transparent to
        drivers and devices.
  \item A bus implementation MUST ensure that any use of \field{token} for
        request/response correlation or to compensate for transport-specific
        reordering remains transparent to drivers and devices.
\end{itemize}

A message originator MUST transmit reserved header bits and unspecified header
values as zero, and receivers MUST ignore those bits on receive to preserve
forward compatibility.

Virtqueue descriptors and buffer data are exchanged via shared memory or other
DMA-accessible memory, as with other virtio transports. Transport messages are
used for control operations, configuration, and notifications, not for
virtqueue data transfer.

\subsubsection{Message Ordering}
\label{sec:Virtio Transport Options / Virtio Over Messages / Basic Concepts / Ordering}

Transport messages fall into two classes: requests (which expect responses) and
events (which are one-way). Drivers and devices do not interpret the
\field{token} field directly. A bus implementation MAY have multiple
outstanding requests for the same device number and MAY use the bus-owned
\field{token} field to correlate each response with its originating request.
The device copies the received request token into the matching response, as
required by
\ref{devicenormative:Virtio Transport Options / Virtio Over Messages / Basic Concepts / Common Message Format / Device}.
The bus may transparently rewrite token values on either leg while preserving
request/response correlation for both endpoints.
Events are asynchronous notifications and MAY be delivered at any time
relative to request/response traffic. Messages for different device numbers MAY
be interleaved by the bus.
Some response payloads repeat request parameters (for example, indices or
offsets). Drivers and devices MAY use these repeated values for validation,
but request/response association is determined by the bus implementation's
correlation state, not by response ordering or by those repeated values.

\busnormative{\paragraph}{Message Ordering}{Virtio Transport Options / Virtio Over Messages / Basic Concepts / Ordering / Bus}
\begin{itemize}
  \item For each transport request that expects a response, a bus
        implementation MUST preserve sufficient correlation state to complete
        that request with exactly one protocol response or one
        transport-visible failure.
  \item A bus implementation MAY forward multiple requests concurrently for the
        same device number.
  \item For each driver/device pair, a bus implementation MUST preserve request
        order when presenting requests to the device-side transport.
  \item For concurrently outstanding requests, a bus implementation MAY deliver
        responses in an order different from the corresponding request order,
        provided it correlates each response with the originating request
        before exposing completion to the driver side.
  \item A bus implementation MUST NOT delay an event solely to preserve a
        request-before-event or response-before-event ordering rule.
\end{itemize}

\devicenormative{\paragraph}{Message Ordering}{Virtio Transport Options / Virtio Over Messages / Basic Concepts / Ordering / Device}
\begin{itemize}
  \item A device MUST process requests in receive order for each driver/device
        pair.
  \item A device MUST send exactly one response for each valid, supported
        request it receives and accepts for processing.
  \item For malformed or unsupported requests, a device MUST discard the
        request without sending a protocol response.
\end{itemize}

\subsection{Device Discovery}\label{sec:Virtio Transport Options / Virtio Over Messages / Device Discovery}

A virtio-msg implementation can learn about devices via bus-level enumeration
(e.g., \busref{GET_DEVICES}), platform data (e.g., device tree, ACPI), or
hypervisor/firmware interfaces. The following informative text describes
typical flows; the normative obligations follow.

Bus implementations discover their devices through a mix of platform data,
hypervisor-provided enumeration, or message-based queries such as
\busref{GET_DEVICES}. This specification defines \busref{GET_DEVICES} for
environments that prefer message-driven enumeration, but it does not require
its use when equivalent information is already known out-of-band.

\busnormative{\paragraph}{Device Discovery}{Virtio Transport Options / Virtio Over Messages / Device Discovery / Bus}
\begin{itemize}
  \item A bus implementation MUST ensure that every device number it exposes to
        the driver has been validated via either platform data or a successful
        query such as \busref{GET_DEVICES}.
  \item Once a device number is deemed present, the bus implementation MUST
        expose that number to the driver side so the driver can issue
        \msgref{GET_DEVICE_INFO}.
  \item If a bus implementation provides an alternative enumeration mechanism
        (e.g., ACPI, device tree), it MAY omit \busref{GET_DEVICES} provided the
        alternative delivers equivalent information to the driver.
\end{itemize}

\subsection{Device Initialization}
\label{sec:Virtio Transport Options / Virtio Over Messages / Device Initialization}

After a bus has identified a virtio-msg device (whether by message-based
enumeration or platform-specific discovery), the driver undertakes a series of
steps to configure and ready the device for normal operation. This section
details the recommended order of operations and the associated messages.

\subsubsection{Initialization Flow Overview}
\label{sec:Virtio Transport Options / Virtio Over Messages / Device Initialization / Overview}

A typical device initialization flow includes:
\begin{enumerate}
  \item \textbf{Early Discovery Exception:}
    Query static parameters with \msgref{GET_DEVICE_INFO}. This is the only
    transport message exception allowed before normal initialization sequencing.
  \item \textbf{Reset And Baseline Status:}
    Reset the device and transition status through ACKNOWLEDGE then DRIVER via
    \msgref{SET_DEVICE_STATUS}.
  \item \textbf{Negotiate Device Features:}
    Read device feature blocks via \msgref{GET_DEVICE_FEATURES}, decide which
    to enable, and send the selected subset via one or more
    \msgref{SET_DRIVER_FEATURES} requests. Verify FEATURES\_OK via
    \msgref{SET_DEVICE_STATUS}. Transport feature bits are bus-instance
    parameters and are determined separately from this per-device exchange.
  \item \textbf{Device-Specific Setup:}
    Read or write configuration data with
    \msgref{GET_CONFIG}/\msgref{SET_CONFIG} using the active configuration
    semantics profile defined in
    \ref{sec:Virtio Transport Options / Virtio Over Messages / Basic Concepts / Configuration Semantics Profiles},
    then configure each virtqueue via \msgref{SET_VQUEUE} and verify settings
    with \msgref{GET_VQUEUE}.
  \item \textbf{Finalize Initialization:}
    Use \msgref{SET_DEVICE_STATUS} to transition to DRIVER\_OK.
\end{enumerate}

After \msgref{GET_DEVICE_INFO}, a driver MUST follow the core virtio status
progression for initialization: reset, ACKNOWLEDGE, DRIVER, feature
negotiation and FEATURES\_OK verification, device-specific setup, then
DRIVER\_OK.

\drivernormative{\paragraph}{Initialization Flow}{Virtio Transport Options / Virtio Over Messages / Device Initialization / Overview / Driver}
\begin{itemize}
  \item A driver MUST issue \msgref{GET_DEVICE_INFO} before attempting feature
        negotiation or queue setup.
  \item A driver MUST complete device feature negotiation via
        \msgref{GET_DEVICE_FEATURES}/\msgref{SET_DRIVER_FEATURES}
        and confirm the FEATURES\_OK state via \msgref{SET_DEVICE_STATUS} before
        enabling virtqueues.
  \item A driver MUST configure each virtqueue via \msgref{SET_VQUEUE} and
        confirm its parameters (e.g., with \msgref{GET_VQUEUE}) before marking
        the queue ready for I/O.
  \item A driver MUST drive the device status transitions using
        \msgref{SET_DEVICE_STATUS}, ensuring the device reaches DRIVER\_OK
        before issuing normal I/O.
\end{itemize}

\subsubsection{Device Information}
\label{sec:Virtio Transport Options / Virtio Over Messages / Device Initialization / Device Information}

Once a device number is confirmed, the driver uses \msgref{GET_DEVICE_INFO} to
retrieve static identification data (device ID, vendor ID), the number of
feature blocks, configuration space size, maximum virtqueues, and any admin
virtqueue details. This information determines which virtio driver should bind
to the device, how many feature blocks to query, and how much configuration
space is valid.

\drivernormative{\paragraph}{Device Information}{Virtio Transport Options / Virtio Over Messages / Device Initialization / Device Information / Driver}
\begin{itemize}
  \item A driver MUST use the configuration size reported via
        \msgref{GET_DEVICE_INFO} to bound any offsets and lengths supplied in
        \msgref{GET_CONFIG} and \msgref{SET_CONFIG} requests.
\end{itemize}

\subsubsection{Feature Negotiation}
\label{sec:Virtio Transport Options / Virtio Over Messages / Device Initialization / Device Features}

Drivers read available device feature bits in 32-bit blocks using
\msgref{GET_DEVICE_FEATURES}; the device returns the requested bitfields,
padding with zeros for any out-of-range blocks. To enable selected device
features, the driver sends the desired subset with one or more
\msgref{SET_DRIVER_FEATURES} requests and then updates the device status via
\msgref{SET_DEVICE_STATUS}, checking whether the FEATURES\_OK bit remains set.
Transport feature bits are bus-instance
parameters and are not part of this message exchange. Zero-padding and
block-scoped update semantics are defined in
\ref{sec:Virtio Transport Options / Virtio Over Messages / Basic Concepts / Device Feature Blocks}.

\drivernormative{\paragraph}{Feature Negotiation}{Virtio Transport Options / Virtio Over Messages / Device Initialization / Device Features / Driver}
\begin{itemize}
  \item A driver MUST use \msgref{GET_DEVICE_FEATURES} to discover the feature
        bits offered by the device and MUST write back only the features it
        intends to enable via \msgref{SET_DRIVER_FEATURES}.
  \item A driver MUST NOT negotiate VIRTIO\_F\_NOTIF\_CONFIG\_DATA.
  \item After writing the desired features, the driver MUST attempt to set the
        FEATURES\_OK bit using \msgref{SET_DEVICE_STATUS} and MUST check the
        returned status to ensure the device accepted the set.
\end{itemize}

\devicenormative{\paragraph}{Feature Negotiation}{Virtio Transport Options / Virtio Over Messages / Device Initialization / Device Features / Device}
\begin{itemize}
  \item A device MUST NOT offer VIRTIO\_F\_NOTIF\_CONFIG\_DATA.
  \item If the device cannot support the driver-selected feature set presented
        before FEATURES\_OK is processed, it MUST clear the FEATURES\_OK bit in
        the status returned by \msgref{SET_DEVICE_STATUS}.
\end{itemize}

\subsubsection{Device Configuration}
\label{sec:Virtio Transport Options / Virtio Over Messages / Device Initialization / Device Configuration}

Drivers use \msgref{GET_CONFIG} to read portions of the configuration space by
supplying \field{offset} and \field{length}; the device returns the requested
data plus the configuration generation field. Writing is performed via
\msgref{SET_CONFIG}, which carries the same \field{offset}/\field{length} plus
generation and new data.
Generation-field handling follows the active profile defined in
\ref{sec:Virtio Transport Options / Virtio Over Messages / Basic Concepts / Configuration Semantics Profiles},
with strict profile defining generation-mismatch rejection for \msgref{SET_CONFIG}.
The message definitions for \msgref{GET_CONFIG} and \msgref{SET_CONFIG}
define request validation and completion behavior, while
\ref{sec:Virtio Transport Options / Virtio Over Messages / Basic Concepts / Configuration Semantics Profiles}
defines generation handling and \msgref{EVENT_CONFIG} semantics throughout
initialization and operation.

Configuration space is intended for device control and configuration data.
When a device can instead expose data through a virtqueue, that mechanism is
preferred; configuration reads and writes are asynchronous in virtio-msg and
may require the driver to wait for a response.

\subsubsection{Virtqueue Configuration}
\label{sec:Virtio Transport Options / Virtio Over Messages / Device Initialization / Virtqueue Configuration}

Drivers query virtqueue parameters with \msgref{GET_VQUEUE}, configure them via
\msgref{SET_VQUEUE}, and optionally reset them using \msgref{RESET_VQUEUE} (if
VIRTIO\_F\_RING\_RESET is negotiated). Each queue is typically configured by
reading its maximum size, provisioning descriptor/available/used buffers, and
then calling \msgref{SET_VQUEUE} with the chosen size, memory addresses, and
queue-state operation in \field{flags}. A \msgref{GET_VQUEUE} response with
\field{max_size}=0 indicates that the requested queue is not available for
configuration. The message definitions for \msgref{GET_VQUEUE},
\msgref{SET_VQUEUE}, and \msgref{RESET_VQUEUE} define the normative
queue-availability, reconfiguration, and reset rules; this subsection
describes the usual sequencing.

\subsubsection{Status Information}
\label{sec:Virtio Transport Options / Virtio Over Messages / Device Initialization / Status Information}

Drivers query the device status via \msgref{GET_DEVICE_STATUS} to observe
progress or detect errors, and they drive the Virtio status transitions via
\msgref{SET_DEVICE_STATUS}. Writing zero to the status field requests a device
reset, which can complete synchronously or asynchronously. The normative
status-reporting and reset-completion rules are defined by
\msgref{GET_DEVICE_STATUS} and \msgref{SET_DEVICE_STATUS}.

\subsubsection{Finalizing Initialization}

After configuring virtqueues and agreeing on features, the driver transitions
the device to DRIVER\_OK (and any other required status bits) via
\msgref{SET_DEVICE_STATUS}. At that point, the device is considered ready for
normal virtqueue I/O.

\drivernormative{\paragraph}{Final Status}{Virtio Transport Options / Virtio Over Messages / Device Initialization / Finalizing Initialization / Driver}
\begin{itemize}
  \item A driver MUST set DRIVER\_OK via \msgref{SET_DEVICE_STATUS} only after
        it has completed feature negotiation and initialized all required
        virtqueues.
  \item A driver MUST NOT queue normal I/O until the device reports a status
        that includes DRIVER\_OK.
  \item A driver MUST NOT supply buffers or send driver notifications for a
        virtqueue until that queue has been configured via \msgref{SET_VQUEUE}.
\end{itemize}

\devicenormative{\paragraph}{Final Status}{Virtio Transport Options / Virtio Over Messages / Device Initialization / Finalizing Initialization / Device}
\begin{itemize}
  \item A device MUST NOT process normal virtqueue I/O until the driver has set
        DRIVER\_OK.
  \item Once DRIVER\_OK is set, the device MUST begin processing virtqueue
        requests subject to the negotiated features and queue configurations.
\end{itemize}

\subsection{Device Operation}
\label{sec:Virtio Transport Options / Virtio Over Messages / Device Operation}

Once a virtio-msg device is fully initialized (see
\ref{sec:Virtio Transport Options / Virtio Over Messages / Device Initialization}),
the driver and device exchange messages to perform I/O and respond to
configuration changes. This section details the primary messaging paths for
\emph{driver notifications}, \emph{device notifications}, and the handling of
runtime resets or shutdown sequences.

\subsubsection{Driver Notifications}
\label{sec:Virtio Transport Options / Virtio Over Messages / Device Operation / Driver Notifications}

Drivers and devices MAY select polling-only operation for virtqueue
availability/completion handling. In polling-only mode, \msgref{EVENT_AVAIL}
and \msgref{EVENT_USED} are not semantic requirements. When endpoints use
event-driven notifications, drivers use \msgref{EVENT_AVAIL} to notify the
device of available buffers. A driver-side bus implementation may forward
\msgref{EVENT_AVAIL} in-band or translate it into an out-of-band notification
mechanism. The receiving-side bus implementation can generate
\msgref{EVENT_AVAIL} locally from that signal. Signaling and coalescing are
specified by the bus binding. Generic bus delivery
rules for runtime notifications are defined later in the transport-message
overview.

\drivernormative{\paragraph}{Driver Notifications}{Virtio Transport Options / Virtio Over Messages / Device Operation / Driver Notifications / Driver}
\begin{itemize}
  \item A driver and device MAY select polling-only operation for virtqueue
        notification/completion handling.
  \item If polling-only mode is selected, a driver MAY omit
        \msgref{EVENT_AVAIL} and rely on polling for progress.
  \item If polling-only mode is not selected, a driver MUST use
        \msgref{EVENT_AVAIL} to notify the device of available buffers.
\end{itemize}

\subsubsection{Device Notifications}
\label{sec:Virtio Transport Options / Virtio Over Messages / Device Operation / Device Notifications}

\msgref{EVENT_CONFIG} and \msgref{EVENT_USED} provide asynchronous notifications
from the device to the driver through the bus layer. Endpoints MAY select
polling-only operation for used-buffer completion; in that mode,
\msgref{EVENT_USED} is not a semantic requirement. A device-side transport
implementation originates event-driven notifications using the standardized
event messages. The local bus implementation may forward those events in-band
or translate them into out-of-band signaling. The receiving-side bus
implementation then forwards the originating event or generates the
corresponding transport event locally. Signaling and virtqueue notification
coalescing are specified by the bus binding. Generic bus delivery rules for
runtime notifications are defined later in the transport-message overview.
\devicenormative{\paragraph}{Device Notifications}{Virtio Transport Options / Virtio Over Messages / Device Operation / Device Notifications / Device}
\begin{itemize}
  \item A device transport layer implementation MUST originate
        \msgref{EVENT_USED} notification occurrences for used buffers on a
        virtqueue when polling-only mode is not selected by the endpoints.
  \item If polling-only mode is selected by the endpoints, a device MAY omit
        \msgref{EVENT_USED} and rely on polling-visible queue state changes for
        completion observation.
\end{itemize}

\drivernormative{\paragraph}{Device Notifications}{Virtio Transport Options / Virtio Over Messages / Device Operation / Device Notifications / Driver}
\begin{itemize}
  \item Upon receiving \msgref{EVENT_CONFIG}, a driver SHOULD update its view of
        the configuration using provided data (or by issuing
        \msgref{GET_CONFIG} if necessary) before resuming normal I/O.
  \item If \msgref{EVENT_CONFIG} provides no data and has
        \field{offset} and \field{length} set to 0, a driver SHOULD treat it as
        a generic configuration-change notification and refresh required
        configuration via \msgref{GET_CONFIG}.
  \item Upon receiving \msgref{EVENT_USED}, a driver MUST examine the indicated
        virtqueue (for example, by reading the used ring) to reclaim completed
        buffers.
  \item If polling-only mode is selected by the endpoints, a driver MAY rely on
        polling to observe used-buffer completion and MUST NOT require a
        transport-visible \msgref{EVENT_USED} semantic occurrence.
\end{itemize}

\subsubsection{Configuration Changes During Operation}
\label{sec:Virtio Transport Options / Virtio Over Messages / Device Operation / Configuration Changes During Operation}

Drivers may update configuration fields at runtime using \msgref{SET_CONFIG}
when features such as device modes or limits need to change. Devices can also
update configuration data autonomously but must signal those changes via
\msgref{EVENT_CONFIG}. The same \msgref{GET_CONFIG}, \msgref{SET_CONFIG}, and
\msgref{EVENT_CONFIG} rules continue to apply after initialization, including
the active configuration semantics profile.

\subsubsection{Virtqueue Changes During Operation}
\label{sec:Virtio Transport Options / Virtio Over Messages / Device Operation / Virtqueue Changes During Operation}

Drivers may provision unused virtqueues later in the device lifetime by issuing
\msgref{SET_VQUEUE}. Queue reconfiguration remains constrained to disabled
queues. If a queue is enabled, the driver first resets that queue via
\msgref{RESET_VQUEUE} when VIRTIO\_F\_RING\_RESET is negotiated. The same
\msgref{GET_VQUEUE}, \msgref{SET_VQUEUE}, and \msgref{RESET_VQUEUE} rules
continue to apply after initialization.

\subsubsection{Device Reset and Shutdown}
\label{sec:Virtio Transport Options / Virtio Over Messages / Device Operation / Reset}

Drivers request a device reset by writing 0 to \field{status} via
\msgref{SET_DEVICE_STATUS}. Reset completion can be synchronous or
asynchronous. A driver confirms completion from a
\msgref{SET_DEVICE_STATUS} response \field{status} of 0 or by polling
\msgref{GET_DEVICE_STATUS} until \field{status} is 0. A device may also signal
that a reset is required by setting DEVICE\_NEEDS\_RESET in its
\field{status} and reporting the updated \field{status} to the driver. After
reset completion, the core virtio post-reset requirements continue to apply.
The normative reset-completion rules are defined by
\msgref{GET_DEVICE_STATUS} and \msgref{SET_DEVICE_STATUS}.

\subsubsection{Hotplug and Removal}
\label{sec:Virtio Transport Options / Virtio Over Messages / Device Operation / Hotplug and Removal}

Support for dynamic device addition and removal is optional. A bus binding
defines which of these operations it supports.

If the bus supports dynamic addition or removal of devices, it can advertise
individual changes with \busref{EVENT_DEVICE} (ADDED or REMOVED), or signal a
change requiring bus-specific handling with NO\_DATA. Some platforms may
instead rely on external signals such as ACPI, device tree updates, or
hypervisor events; virtio-msg does not mandate a specific mechanism.

Device addition and removal are managed by the bus implementation. An addition
indication reports that a device is accessible when announced; the device can
subsequently become unavailable at any time. Bus bindings define removal
sequencing, synchronization with driver initialization and operation, and the
conditions for releasing associated resources. A bus binding may provide a
coordinated removal mechanism.

\busnormative{\paragraph}{Hotplug and Removal}{Virtio Transport Options / Virtio Over Messages / Device Operation / Hotplug and Removal / Bus}
\begin{itemize}
  \item A bus implementation that supports dynamic device addition or removal
        and uses \busref{EVENT_DEVICE} to report those changes MUST report each
        supported change with the corresponding ADDED or REMOVED event, or use
        NO\_DATA to invoke bus-specific handling of those changes in place of
        one or more individual events.
  \item Regardless of how hotplug information is delivered, once a new device is
        reported the bus implementation MUST expose the device number to the
        driver side so the driver can issue \msgref{GET_DEVICE_INFO}. When a
        device is removed, the bus implementation MUST notify the driver side.
  \item A bus implementation that supports dynamic device addition or removal
        MUST ensure that removal is handled safely during device initialization
        and operation.
  \item A bus implementation MUST complete transport requests outstanding when
        a device is removed according to
        \ref{busnormative:Virtio Transport Options / Virtio Over Messages / Basic Concepts / Error Signaling / Bus}.
  \item A bus implementation that supports dynamic device addition or removal
        MUST define how the driver side establishes that the removed device can
        no longer access associated resources before those resources are
        released.
\end{itemize}

\subsection{Transport Messages}\label{sec:Virtio Transport Options / Virtio Over Messages / Transport Messages}

Transport messages configure and operate virtio devices once they have been
discovered. Unlike bus messages (which cover enumeration or hotplug), transport
messages focus on feature negotiation, configuration access, virtqueue setup,
and runtime notifications. The subsections below describe each message and how
it is used in the virtio-msg transport.

\subsubsection{Overview}
\label{sec:Virtio Transport Options / Virtio Over Messages / Transport Messages / Overview}

Drivers send transport messages to a specific device number, and the virtio-msg
device replies or emits events accordingly. The bus simply forwards these
messages after enforcing generic checks such as maximum size or verifying that
the target device exists. Most transport messages are request/response pairs,
with events reserved for asynchronous notifications.

\paragraph{Message IDs and initiating side}

In the table below, the initiating side identifies which side of the
transport-message exchange starts the message.

\begin{tabular}{|l|l|l|}
\hline
Name & ID & Initiating side \\
\hline
\hline
Reserved                   & 0x0 & \\
\hline
Reserved                   & 0x1 & \\
\hline
\msgref{GET_DEVICE_INFO}   & 0x2  & Driver side \\
\hline
\msgref{GET_DEVICE_FEATURES} & 0x3  & Driver side \\
\hline
\msgref{SET_DRIVER_FEATURES} & 0x4  & Driver side \\
\hline
\msgref{GET_CONFIG}        & 0x5  & Driver side \\
\hline
\msgref{SET_CONFIG}        & 0x6  & Driver side \\
\hline
\msgref{GET_DEVICE_STATUS} & 0x7  & Driver side \\
\hline
\msgref{SET_DEVICE_STATUS} & 0x8  & Driver side \\
\hline
\msgref{GET_VQUEUE}        & 0x9  & Driver side \\
\hline
\msgref{SET_VQUEUE}        & 0xA  & Driver side \\
\hline
\msgref{RESET_VQUEUE}      & 0xB  & Driver side \\
\hline
\msgref{GET_SHM}           & 0xC  & Driver side \\
\hline
\msgref{EVENT_CONFIG}      & 0x40 & Device side \\
\hline
\msgref{EVENT_AVAIL}       & 0x41 & Driver side \\
\hline
\msgref{EVENT_USED}        & 0x42 & Device side \\
\hline
\end{tabular}

Transport-message ID bit encoding is defined in
\ref{sec:Virtio Transport Options / Virtio Over Messages / Basic Concepts / Common Message Format}.

\devicenormative{\paragraph}{Core Messages}{Virtio Transport Options / Virtio Over Messages / Transport Messages / Mandatory / Device}
\begin{itemize}
  \item A device MUST implement \msgref{GET_DEVICE_INFO},
        \msgref{GET_DEVICE_FEATURES}, \msgref{SET_DRIVER_FEATURES},
        \msgref{GET_CONFIG}, \msgref{SET_CONFIG}, \msgref{GET_DEVICE_STATUS},
        \msgref{SET_DEVICE_STATUS}, \msgref{GET_VQUEUE}, \msgref{SET_VQUEUE},
        \msgref{RESET_VQUEUE}, and \msgref{GET_SHM}.
\end{itemize}

\drivernormative{\paragraph}{Core Messages}{Virtio Transport Options / Virtio Over Messages / Transport Messages / Mandatory / Driver}
\begin{itemize}
  \item A driver MUST support sending and interpreting
        \msgref{GET_DEVICE_INFO}, \msgref{GET_DEVICE_FEATURES},
        \msgref{SET_DRIVER_FEATURES}, \msgref{GET_CONFIG},
        \msgref{SET_CONFIG}, \msgref{GET_DEVICE_STATUS},
        \msgref{SET_DEVICE_STATUS}, \msgref{GET_VQUEUE},
        \msgref{SET_VQUEUE}, and \msgref{RESET_VQUEUE} in order to
        initialize and operate a virtio-msg device.
\end{itemize}

As described in Sections
\ref{sec:Virtio Transport Options / Virtio Over Messages / Device Operation / Driver Notifications}
and
\ref{sec:Virtio Transport Options / Virtio Over Messages / Device Operation / Device Notifications},
runtime notifications (\msgref{EVENT_AVAIL}, \msgref{EVENT_USED}, and
\msgref{EVENT_CONFIG}) are defined as transport messages. A bus binding can
carry these notifications as messages or use bus-specific signaling, such as
hardware or software interrupts, from which the receiving-side bus
implementation generates transport events locally. An event message need not
be sent over the bus. There need not be a one-to-one correspondence between
bus signals and transport events. Bus bindings define the signaling mechanism
and coalescing policy for virtqueue notifications.

\busnormative{\paragraph}{Runtime Notifications}{Virtio Transport Options / Virtio Over Messages / Transport Messages / Runtime Notifications / Bus}
\begin{itemize}
  \item A bus implementation MUST define how runtime notifications
        (\msgref{EVENT_AVAIL}, \msgref{EVENT_USED}, and \msgref{EVENT_CONFIG})
        are delivered to the receiving transport, including any bus-specific
        signaling, local event synthesis, coalescing, and required
        configuration.
  \item A bus implementation MAY coalesce \msgref{EVENT_AVAIL} or
        \msgref{EVENT_USED} notifications for the same virtqueue and direction,
        provided it preserves the required virtqueue notifications and progress,
        including any negotiated notification-data semantics.
  \item A bus implementation MUST deliver each \msgref{EVENT_CONFIG}
        notification as a forwarded or locally synthesized transport event.
  \item A bus implementation MUST NOT expose both forwarded and synthesized
        transport events for the same notification, or expose a bus signal as
        an additional transport event when the notification has already been
        delivered by forwarding or synthesis.
  \item If polling-only mode is selected by the endpoints for availability
        handling or used-buffer completion, any bus-side polling, forwarding,
        or synthesis logic for that function MUST remain transparent and MUST
        NOT create a transport-visible \msgref{EVENT_AVAIL} or
        \msgref{EVENT_USED} semantic occurrence.
  \item A bus implementation MUST use bounded timeout and polling policies for
        runtime-notification delivery. If a required transport-visible event
        cannot be delivered within the bounded policy, it MUST surface a
        transport-visible failure and MUST define recovery behavior using the
        reset and status mechanisms.
\end{itemize}

\msgdef{GET_DEVICE_INFO}

The driver issues \msgref{GET_DEVICE_INFO} to discover static identification
data and limits for a device. The response includes \field{device_uuid}, a
16-byte UUID value. The nil UUID defined by
\hyperref[intro:rfc4122]{[RFC4122]} indicates that no UUID is provided.
\field{admin_vq_start} and \field{admin_vq_count} identify a contiguous range
of administration virtqueues within \field{max_virtqueues}. If
\field{admin_vq_count} is 0, the device exposes no administration virtqueues
and \field{admin_vq_start} is 0.
\msgref{GET_DEVICE_INFO} is the only transport message that may be issued as an
early-discovery exception before the normal initialization sequence begins.
After \msgref{GET_DEVICE_INFO}, drivers follow the standard virtio
initialization sequence and use other transport messages only within that flow.

\begin{lstlisting}
struct virtio_msg_get_device_info_req {
        /* no payload */
};

struct virtio_msg_get_device_info_resp {
        le32 device_id;        /* virtio device type */
        le32 vendor_id;        /* implementation-defined vendor ID */
        u8   device_uuid[16];  /* UUID, or nil UUID if unavailable */
        le32 num_feature_blocks; /* total number of 32-bit feature blocks */
        le32 config_size;      /* bytes in the device configuration space */
        le32 max_virtqueues;   /* maximum virtqueues supported */
        le32 admin_vq_start;   /* index of the first admin virtqueue */
        le32 admin_vq_count;   /* number of admin virtqueues */
};
\end{lstlisting}

\devicenormative{\paragraph}{GET\_DEVICE\_INFO}{Virtio Transport Options / Virtio Over Messages / Transport Messages / VIRTIO_MSG_GET_DEVICE_INFO / Device}
\begin{itemize}
  \item A device MUST respond to \msgref{GET_DEVICE_INFO} while inactive and
        MUST supply consistent values for device ID, vendor ID, feature-block
        count, configuration size, and virtqueue limits.
  \item A device MUST report \field{num_feature_blocks} large enough to cover
        every feature bit the device offers.
  \item A device MUST report \field{max_virtqueues} no greater than 65536.
  \item \field{max_virtqueues} MUST include any administration virtqueues
        described by \field{admin_vq_start} and \field{admin_vq_count}.
  \item If \field{admin_vq_count} is 0, a device MUST set
        \field{admin_vq_start} to 0.
  \item If \field{admin_vq_count} is greater than 0, a device MUST ensure
        \field{admin_vq_start} + \field{admin_vq_count} does not exceed
        \field{max_virtqueues}.
  \item A device MAY set \field{device_uuid} to a UUID value that can be used
        as a unique identifier.
  \item If a device sets \field{device_uuid} to a non-nil value, it SHOULD use
        a version 4 UUID as specified by
        \hyperref[intro:rfc4122]{[RFC4122]}.
  \item If a device does not support UUID reporting or does not have an
        identifier, it MUST set \field{device_uuid} to the nil UUID.
\end{itemize}

\drivernormative{\paragraph}{GET\_DEVICE\_INFO}{Virtio Transport Options / Virtio Over Messages / Transport Messages / VIRTIO_MSG_GET_DEVICE_INFO / Driver}
\begin{itemize}
  \item A driver SHOULD treat the response as the authoritative source for
        feature-block count, configuration size, and virtqueue limits before
        proceeding with initialization.
  \item A driver MUST interpret \field{admin_vq_count}=0 as indicating that the
        device exposes no administration virtqueues.
  \item If \field{admin_vq_count} is greater than 0, a driver MUST treat the
        range starting at \field{admin_vq_start} and spanning
        \field{admin_vq_count} entries as part of \field{max_virtqueues}.
  \item A driver MUST interpret a nil UUID in \field{device_uuid} as indicating
        that no UUID is provided.
\end{itemize}

\msgdef{GET_DEVICE_FEATURES}

Drivers retrieve device feature bits in 32-bit blocks via
\msgref{GET_DEVICE_FEATURES}; the response echoes the requested block index and
returns one or more 32-bit values with the device-offered feature bits in that
range.

\begin{lstlisting}
struct virtio_msg_get_device_features_req {
        le32 block_index; /* starting block (0 == bits 0-31) */
        le32 num_blocks;  /* number of 32-bit blocks requested */
};

struct virtio_msg_get_device_features_resp {
        le32 block_index; /* echoed starting block */
        le32 num_blocks;  /* echoed number of blocks */
        le32 features[];  /* num_blocks entries, zero-padded if needed */
};
\end{lstlisting}

Zero-padding for out-of-range blocks is defined in
\ref{sec:Virtio Transport Options / Virtio Over Messages / Basic Concepts / Device Feature Blocks}.

\drivernormative{\paragraph}{GET\_DEVICE\_FEATURES}{Virtio Transport Options / Virtio Over Messages / Transport Messages / VIRTIO_MSG_GET_DEVICE_FEATURES / Driver}
\begin{itemize}
  \item A driver MUST treat the returned data as 32-bit feature blocks starting
        at the requested index, padding with zeros as needed, before deciding
        which features to enable.
\end{itemize}

\devicenormative{\paragraph}{GET\_DEVICE\_FEATURES}{Virtio Transport Options / Virtio Over Messages / Transport Messages / VIRTIO_MSG_GET_DEVICE_FEATURES / Device}
\begin{itemize}
  \item A device MUST echo the requested \field{block_index} and
        \field{num_blocks} in its response.
\end{itemize}

\msgdef{SET_DRIVER_FEATURES}

Drivers use \msgref{SET_DRIVER_FEATURES} to convey the subset of feature
bits they wish to enable. The payload mirrors \msgref{GET_DEVICE_FEATURES},
supplying an index, count, and 32-bit feature data for the selected blocks.
Each \msgref{SET_DRIVER_FEATURES} request updates only the addressed blocks;
blocks outside the addressed range are unchanged. Drivers MAY send multiple
\msgref{SET_DRIVER_FEATURES} requests at different block offsets before setting
FEATURES\_OK. The response has no payload.

\begin{lstlisting}
struct virtio_msg_set_driver_features_req {
        le32 block_index; /* starting block being selected */
        le32 num_blocks;  /* number of blocks provided */
        le32 features[];  /* num_blocks entries with desired bits */
};

struct virtio_msg_set_driver_features_resp {
        /* no payload */
};
\end{lstlisting}

\drivernormative{\paragraph}{SET\_DRIVER\_FEATURES}{Virtio Transport Options / Virtio Over Messages / Transport Messages / VIRTIO_MSG_SET_DRIVER_FEATURES / Driver}
\begin{itemize}
  \item A driver MUST set only the feature bits it intends to enable in the
        blocks supplied to \msgref{SET_DRIVER_FEATURES}.
  \item A driver MAY use multiple \msgref{SET_DRIVER_FEATURES} requests with
        different \field{block_index} values to build the final driver-selected
        feature set before setting FEATURES\_OK.
\end{itemize}

Block-scoped update semantics for \msgref{SET_DRIVER_FEATURES} are defined in
\ref{sec:Virtio Transport Options / Virtio Over Messages / Basic Concepts / Device Feature Blocks}.

\msgdef{GET_CONFIG}

\msgref{GET_CONFIG} reads a portion of the configuration space. The request
specifies the byte \field{offset} and \field{length}; the response echoes those
values, returns the data, and includes the current generation count. Generation
semantics follow the active configuration profile defined in
\ref{sec:Virtio Transport Options / Virtio Over Messages / Basic Concepts / Configuration Semantics Profiles}.

\begin{lstlisting}
struct virtio_msg_get_config_req {
        le32 offset; /* byte offset into configuration space */
        le32 length; /* number of bytes requested */
};

struct virtio_msg_get_config_resp {
        le32 generation; /* current generation count */
        le32 offset;     /* echoed offset */
        le32 length;     /* echoed length */
        u8   data[];     /* configuration payload */
};
\end{lstlisting}

\drivernormative{\paragraph}{GET\_CONFIG}{Virtio Transport Options / Virtio Over Messages / Transport Messages / VIRTIO_MSG_GET_CONFIG / Driver}
\begin{itemize}
  \item A driver MUST ensure the requested \field{offset} and \field{length}
        in \msgref{GET_CONFIG} stay within the configuration size reported by
        \msgref{GET_DEVICE_INFO}.
\end{itemize}

\devicenormative{\paragraph}{GET\_CONFIG}{Virtio Transport Options / Virtio Over Messages / Transport Messages / VIRTIO_MSG_GET_CONFIG / Device}
\begin{itemize}
  \item A device MUST echo the requested \field{offset} and \field{length}
        in its response.
  \item A device MUST include its current generation count in the
        \field{generation} field of every \msgref{GET_CONFIG} response.
\end{itemize}

\msgdef{SET_CONFIG}

\msgref{SET_CONFIG} writes a portion of configuration space, supplying the
\field{offset}, \field{length}, the driver's view of the generation count, and
the new data. The device echoes the \field{offset}, returns the
resulting generation count, and reports in response \field{length} whether the
requested nonzero write was applied in full. A request with \field{length}=0 is
a valid no-op. For a request with \field{length}>0, the device either applies
all requested bytes and returns response \field{length} equal to the request
\field{length}, or applies nothing and returns response \field{length}=0 with no
payload bytes.
Generation-mismatch rejection is defined only for strict profile.
Generation semantics follow the active configuration profile defined in
\ref{sec:Virtio Transport Options / Virtio Over Messages / Basic Concepts / Configuration Semantics Profiles}.

\begin{lstlisting}
struct virtio_msg_set_config_req {
        le32 generation; /* driver's view; used for mismatch rejection only in
                          * strict profile (MUST be 0 otherwise) */
        le32 offset;     /* byte offset being written */
        le32 length;     /* number of bytes requested; 0 is a valid no-op */
        u8   data[];     /* configuration payload */
};

struct virtio_msg_set_config_resp {
        le32 generation; /* resulting generation count */
        le32 offset;     /* echoed offset */
        le32 length;     /* request length if fully applied; 0 if not applied */
        u8   data[];     /* echoed data for a fully applied nonzero write; omitted if length is 0 */
};
\end{lstlisting}

\drivernormative{\paragraph}{SET\_CONFIG}{Virtio Transport Options / Virtio Over Messages / Transport Messages / VIRTIO_MSG_SET_CONFIG / Driver}
\begin{itemize}
  \item A driver MUST keep the \field{offset} and \field{length} within the
        reported configuration size in each \msgref{SET_CONFIG} request.
  \item A driver MAY issue \msgref{SET_CONFIG} with \field{length}=0 as a
        no-op.
  \item If a driver issues \msgref{SET_CONFIG} with request \field{length}>0,
        it MUST treat response \field{length} equal to the request
        \field{length} as fully applied and response \field{length}=0 as not
        applied; partial nonzero lengths are not valid success indications.
\end{itemize}

\devicenormative{\paragraph}{SET\_CONFIG}{Virtio Transport Options / Virtio Over Messages / Transport Messages / VIRTIO_MSG_SET_CONFIG / Device}
\begin{itemize}
  \item A device MUST echo the requested \field{offset} in its response.
  \item If a device sets response \field{length} to a nonzero value, it MUST
        echo the request \field{data} in its response.
  \item A device MUST set response \field{generation} to its current generation
        count in every \msgref{SET_CONFIG} response.
  \item If a device accepts \msgref{SET_CONFIG} with request \field{length}=0,
        it MUST treat the request as a no-op.
  \item For \msgref{SET_CONFIG} requests with request \field{length}>0, a
        device MUST either apply all requested bytes and set response
        \field{length} to the request \field{length}, or apply no bytes and set
        response \field{length} to 0; partial nonzero application is not
        allowed.
  \item If a device sets response \field{length} to 0, it MUST NOT include
        response \field{data} bytes.
\end{itemize}

\msgdef{GET_DEVICE_STATUS}

\msgref{GET_DEVICE_STATUS} reads the current device status (e.g., ACKNOWLEDGE,
DRIVER, DRIVER\_OK, FEATURES\_OK, or DEVICE\_NEEDS\_RESET bits). The request has
no payload; the response returns the 32-bit status value.

\begin{lstlisting}
struct virtio_msg_get_device_status_req {
        /* no payload */
};

struct virtio_msg_get_device_status_resp {
        le32 status; /* current device status bits */
};
\end{lstlisting}

\drivernormative{\paragraph}{GET\_DEVICE\_STATUS}{Virtio Transport Options / Virtio Over Messages / Transport Messages / VIRTIO_MSG_GET_DEVICE_STATUS / Driver}
\begin{itemize}
  \item A driver MAY send \msgref{GET_DEVICE_STATUS} whenever it needs the
        current device status bits.
  \item If device status reports DEVICE\_NEEDS\_RESET, a driver SHOULD write 0
        via \msgref{SET_DEVICE_STATUS} and reinitialize the device before
        resuming I/O.
  \item When waiting for reset completion after writing status 0 via
        \msgref{SET_DEVICE_STATUS} and synchronous completion is not indicated,
        a driver MUST poll \msgref{GET_DEVICE_STATUS} until
        \field{status}=0 before continuing initialization.
\end{itemize}

\devicenormative{\paragraph}{GET\_DEVICE\_STATUS}{Virtio Transport Options / Virtio Over Messages / Transport Messages / VIRTIO_MSG_GET_DEVICE_STATUS / Device}
\begin{itemize}
  \item A device MUST report its current status accurately via
        \msgref{GET_DEVICE_STATUS}, including whether FEATURES\_OK or
        DEVICE\_NEEDS\_RESET are set.
  \item If a device encounters an unrecoverable error that requires driver
        intervention, it SHOULD set DEVICE\_NEEDS\_RESET in the status it
        reports to the driver.
\end{itemize}

\msgdef{SET_DEVICE_STATUS}

\msgref{SET_DEVICE_STATUS} writes a new device status value. Drivers use it to
progress through the virtio-defined states or to request a reset by writing 0.
The response reports the \field{status} value observed when the response is
generated. For a reset request (\field{status} == 0), a response
\field{status} of 0 indicates reset completion; a non-zero \field{status}
indicates that reset completion may still be pending. The resulting
\field{status} may differ from the requested \field{status} (for example, if
the device refuses FEATURES\_OK or sets DEVICE\_NEEDS\_RESET). After reset
completion, core virtio post-reset requirements continue to apply.

\begin{lstlisting}
struct virtio_msg_set_device_status_req {
        le32 status; /* desired device status value */
};

struct virtio_msg_set_device_status_resp {
        le32 status; /* resulting device status */
};
\end{lstlisting}

\drivernormative{\paragraph}{SET\_DEVICE\_STATUS}{Virtio Transport Options / Virtio Over Messages / Transport Messages / VIRTIO_MSG_SET_DEVICE_STATUS / Driver}
\begin{itemize}
  \item A driver MUST write 0 via \msgref{SET_DEVICE_STATUS} to request a device
        reset.
  \item A driver that needs to confirm reset completion MUST treat a
        \msgref{SET_DEVICE_STATUS} response \field{status} of 0 as completion;
        otherwise it MUST poll \msgref{GET_DEVICE_STATUS} until the device
        reports \field{status}=0 before reinitializing the device.
\end{itemize}

\devicenormative{\paragraph}{SET\_DEVICE\_STATUS}{Virtio Transport Options / Virtio Over Messages / Transport Messages / VIRTIO_MSG_SET_DEVICE_STATUS / Device}
\begin{itemize}
  \item A device MAY clear FEATURES\_OK or set DEVICE\_NEEDS\_RESET in its
        response if it cannot accept the requested \field{status}, but it MUST
        report the resulting \field{status} accurately.
  \item Upon receiving \msgref{SET_DEVICE_STATUS} with \field{status} set to 0,
        a device MUST start reset processing and MUST eventually complete reset
        by reporting \field{status}=0.
  \item A device MAY complete reset before sending the
        \msgref{SET_DEVICE_STATUS} response (synchronous completion) or after
        sending it (asynchronous completion).
  \item A device MUST NOT report \field{status}=0 until reset completion.
\end{itemize}

\msgdef{GET_VQUEUE}

\msgref{GET_VQUEUE} returns information about a specific virtqueue, including
its maximum size, current size, enabled state, and, if already configured, the
descriptor, driver, and device area addresses, as defined in
\ref{sec:Virtio Transport Options / Virtio Over Messages / Basic Concepts / Memory Addresses}.
For unavailable,
unimplemented, or out-of-range queue indexes, the response uses a
deterministic zero sentinel while still echoing the requested \field{index}.

\begin{lstlisting}
struct virtio_msg_get_vqueue_req {
        le32 index; /* virtqueue index */
};

struct virtio_msg_get_vqueue_resp {
        le32 index;        /* echoed virtqueue index */
        le32 max_size;     /* maximum queue size (0 if unavailable) */
        le32 cur_size;     /* current size (0 if unconfigured/unavailable) */
        le32 flags;        /* bit0: enabled state, 0 if unavailable */
        le64 desc_addr;    /* descriptor area address (0 if unavailable) */
        le64 driver_addr;  /* driver area address (0 if unavailable) */
        le64 device_addr;  /* device area address (0 if unavailable) */
};
\end{lstlisting}

\drivernormative{\paragraph}{GET\_VQUEUE}{Virtio Transport Options / Virtio Over Messages / Transport Messages / VIRTIO_MSG_GET_VQUEUE / Driver}
\begin{itemize}
  \item Before configuring a virtqueue, a driver MUST query that queue with
        \msgref{GET_VQUEUE} and MUST treat \field{max_size}=0 as unavailable,
        unimplemented, or out of range.
  \item When issuing \msgref{SET_VQUEUE} for a queue with nonzero
        \field{max_size}, a driver MUST NOT request a queue size greater than
        the \field{max_size} reported by \msgref{GET_VQUEUE}.
\end{itemize}

\devicenormative{\paragraph}{GET\_VQUEUE}{Virtio Transport Options / Virtio Over Messages / Transport Messages / VIRTIO_MSG_GET_VQUEUE / Device}
\begin{itemize}
  \item A device MUST echo the requested \field{index} in its response.
  \item For an available and implemented virtqueue index, a device MUST report
        accurate maxima and current queue sizes and MUST return zero as the
        current size if the queue has not yet been configured.
  \item For an unavailable, unimplemented, or out-of-range virtqueue index
        (\field{index} >= \field{max_virtqueues} from \msgref{GET_DEVICE_INFO}),
        a device MUST return the deterministic unavailable-queue sentinel:
        \field{max_size}=0, \field{cur_size}=0, \field{flags}=0,
        \field{desc_addr}=0, \field{driver_addr}=0, and \field{device_addr}=0,
        while echoing the requested \field{index}.
  \item A device MUST report queue enabled state in \field{flags} bit 0 and
        MUST set \field{flags} bits 31:1 to zero.
\end{itemize}

\msgdef{SET_VQUEUE}

\msgref{SET_VQUEUE} updates a virtqueue's parameters and queue state. Address
fields follow
\ref{sec:Virtio Transport Options / Virtio Over Messages / Basic Concepts / Memory Addresses}.
The driver selects a queue index and uses \field{flags} to indicate whether queue
state should be enabled, kept disabled, or kept unchanged, and which queue
fields should be ignored versus updated.

\begin{lstlisting}
struct virtio_msg_set_vqueue_req {
        le32 index;       /* virtqueue index */
        le32 flags;       /* state op + ignore bits */
        le32 size;        /* number of descriptors */
        le32 reserved;    /* must be zero */
        le64 desc_addr;   /* descriptor area address */
        le64 driver_addr; /* driver area address */
        le64 device_addr; /* device area address */
};

struct virtio_msg_set_vqueue_resp {
        /* no payload */
};
\end{lstlisting}

\field{flags} is defined as:
\begin{itemize}
  \item Bits [1:0] (\textbf{State Operation}):
    \begin{itemize}
      \item 0: keep queue disabled.
      \item 1: enable queue.
      \item 2: keep current queue state.
      \item 3: reserved.
    \end{itemize}
  \item Bit 2 (\textbf{size\_ignore}): 0=apply \field{size}, 1=ignore.
  \item Bit 3 (\textbf{desc\_addr\_ignore}): 0=apply \field{desc_addr},
        1=ignore.
  \item Bit 4 (\textbf{driver\_addr\_ignore}): 0=apply \field{driver_addr},
        1=ignore.
  \item Bit 5 (\textbf{device\_addr\_ignore}): 0=apply \field{device_addr},
        1=ignore.
  \item Bits [31:6]: reserved, MUST be zero.
\end{itemize}

When \field{flags} is zero, \msgref{SET_VQUEUE} uses current behavior:
if the queue is disabled, all queue fields in the request are applied and the
queue remains disabled.

\drivernormative{\paragraph}{SET\_VQUEUE}{Virtio Transport Options / Virtio Over Messages / Transport Messages / VIRTIO_MSG_SET_VQUEUE / Driver}
\begin{itemize}
  \item A driver MUST set \field{reserved} to zero and MUST set \field{flags}
        bits [31:6] to zero.
  \item A driver MUST NOT use state operation value 3 in \field{flags} bits
        [1:0].
  \item A driver SHOULD use \msgref{GET_VQUEUE} to determine whether a queue is
        available before relying on \msgref{SET_VQUEUE} to modify it.
  \item A driver MUST NOT use state operation value 0 as a request to disable an
        already enabled queue.
  \item If \field{size\_ignore} is 0, a driver MUST set \field{size} to a value
        not exceeding the maximum reported by \msgref{GET_VQUEUE}.
  \item If any address ignore bit is 0, a driver MUST ensure the corresponding
        address identifies memory with the alignment required by the virtqueue
        format.
  \item A driver MUST NOT attempt to update queue size or queue addresses via
        \msgref{SET_VQUEUE} while the queue is enabled.
  \item If a driver needs to reconfigure an enabled queue, it MUST issue
        \msgref{RESET_VQUEUE} first when VIRTIO\_F\_RING\_RESET is negotiated.
  \item Because \msgref{SET_VQUEUE} has no response payload, a driver that
        needs confirmation of acceptance MUST verify the resulting queue
        parameters and state via \msgref{GET_VQUEUE} before relying on the
        requested updates.
\end{itemize}

\devicenormative{\paragraph}{SET\_VQUEUE}{Virtio Transport Options / Virtio Over Messages / Transport Messages / VIRTIO_MSG_SET_VQUEUE / Device}
\begin{itemize}
  \item For an unavailable, unimplemented, or out-of-range virtqueue index
        (\field{index} >= \field{max_virtqueues} from \msgref{GET_DEVICE_INFO}),
        a device MUST treat \msgref{SET_VQUEUE} as a no-op and MUST NOT modify
        any queue state or parameters.
  \item If \field{reserved} is nonzero, if \field{flags} bits [31:6] are
        nonzero, or if \field{flags} bits [1:0] use state operation value 3, a
        device MUST do nothing (no queue state or parameter change) and MUST
        return the normal empty \msgref{SET_VQUEUE} response.
  \item If \field{flags} requests state operation value 0 for a queue that is
        currently enabled, a device MUST do nothing (no queue state or
        parameter change) and MUST return the normal empty
        \msgref{SET_VQUEUE} response.
  \item A device MUST process \msgref{SET_VQUEUE} updates atomically: it MUST
        either apply all requested field updates and queue-state changes, or
        do nothing (no queue state or parameter change) and return the normal
        empty \msgref{SET_VQUEUE} response.
  \item If an ignore bit is 0, a device MUST apply the corresponding field from
        \msgref{SET_VQUEUE}; if an ignore bit is 1, a device MUST keep the
        current value for that field unchanged.
  \item For a queue that is enabled, if \msgref{SET_VQUEUE} attempts to update
        \field{size}, \field{desc_addr}, \field{driver_addr}, or
        \field{device_addr}, a device MUST do nothing (no queue state or
        parameter change) and MUST return the normal empty
        \msgref{SET_VQUEUE} response.
  \item A device MUST apply queue state according to \field{flags} bits [1:0]
        after processing requested field updates, and MUST NOT perform an
        enabled-to-disabled transition via \msgref{SET_VQUEUE}.
\end{itemize}

\msgdef{RESET_VQUEUE}

\msgref{RESET_VQUEUE} resets an individual virtqueue (if VIRTIO\_F\_RING\_RESET
has been negotiated). It requires only the queue index; the device responds
after the queue has been quiesced. After reset completion, the core virtio
post-reset requirements for that queue continue to apply.

\begin{lstlisting}
struct virtio_msg_reset_vqueue_req {
        le32 index; /* virtqueue index */
};

struct virtio_msg_reset_vqueue_resp {
        /* no payload */
};
\end{lstlisting}

\drivernormative{\paragraph}{RESET\_VQUEUE}{Virtio Transport Options / Virtio Over Messages / Transport Messages / VIRTIO_MSG_RESET_VQUEUE / Driver}
\begin{itemize}
  \item A driver MUST issue \msgref{RESET_VQUEUE} only if
        VIRTIO\_F\_RING\_RESET has been negotiated and it needs to reconfigure
        or recover a specific queue.
  \item A driver SHOULD use \msgref{GET_VQUEUE} to determine whether a queue is
        available before relying on \msgref{RESET_VQUEUE} to affect it.
  \item If an enabled queue requires reconfiguration and
        VIRTIO\_F\_RING\_RESET is negotiated, a driver MUST wait for
        \msgref{RESET_VQUEUE} completion before issuing \msgref{SET_VQUEUE}
        updates for that queue.
\end{itemize}

\devicenormative{\paragraph}{RESET\_VQUEUE}{Virtio Transport Options / Virtio Over Messages / Transport Messages / VIRTIO_MSG_RESET_VQUEUE / Device}
\begin{itemize}
  \item For an unavailable, unimplemented, or out-of-range virtqueue index
        (\field{index} >= \field{max_virtqueues} from \msgref{GET_DEVICE_INFO}),
        a device MUST treat \msgref{RESET_VQUEUE} as a no-op and MUST NOT
        modify any queue state.
  \item If VIRTIO\_F\_RING\_RESET has not been negotiated, a device MUST do
        nothing (no queue state change) and MUST return the normal empty
        \msgref{RESET_VQUEUE} response.
  \item Upon receiving \msgref{RESET_VQUEUE}, a device MUST stop processing the
        indicated queue, reset its internal state for that queue, and allow the
        driver to reconfigure it via \msgref{SET_VQUEUE} while the queue is
        disabled.
\end{itemize}

\msgdef{GET_SHM}

\msgref{GET_SHM} returns the location and size of a device-owned shared memory
region. The request carries the shared-memory region identifier
\field{shmid}; the response echoes \field{shmid} and provides the base
\field{address} and \field{length}. The address follows
\ref{sec:Virtio Transport Options / Virtio Over Messages / Basic Concepts / Memory Addresses}.
A \field{length} of zero
indicates that the specified \field{shmid} is not supported by the device.

\begin{lstlisting}
struct virtio_msg_get_shm_req {
        le32 shmid; /* shared memory region identifier */
};

struct virtio_msg_get_shm_resp {
        le32 shmid;     /* echoed region identifier */
        le32 reserved;  /* must be zero */
        le64 length;    /* region length (0 if nonexistent) */
        le64 address;   /* region address */
};
\end{lstlisting}

\devicenormative{\paragraph}{GET\_SHM}{Virtio Transport Options / Virtio Over Messages / Transport Messages / VIRTIO_MSG_GET_SHM / Device}
\begin{itemize}
  \item A device MUST echo the requested \field{shmid} in its response,
        including when the requested region is not supported.
  \item A device MUST return zero \field{length} if the requested
        \field{shmid} is not supported and MUST report accurate
        \field{address}/\field{length} information for regions it supports.
  \item A device MUST set the \field{reserved} field in \msgref{GET_SHM}
        responses to zero.
\end{itemize}

\msgdef{EVENT_CONFIG}

\msgref{EVENT_CONFIG} notifies the driver about configuration changes or
device-originated asynchronous status changes that require asynchronous
reporting. The payload includes current device status, generation, and
configuration \field{offset}/\field{length}. The configuration data can be
omitted; if omitted, the driver re-fetches via \msgref{GET_CONFIG}. Generation
and \field{offset}/\field{length} semantics are defined in
\ref{sec:Virtio Transport Options / Virtio Over Messages / Basic Concepts / Configuration Semantics Profiles}.
Ordinary driver-originated status writes via \msgref{SET_DEVICE_STATUS} are
already observed synchronously and do not by themselves trigger
\msgref{EVENT_CONFIG}; an event may be sent later only when a distinct
device-originated asynchronous status change occurs afterward.
If the event reports only a status change with no configuration update, the
device sets \field{offset} and \field{length} to zero.

\begin{lstlisting}
struct virtio_msg_event_config {
        le32 device_status; /* new device status value */
        le32 generation;    /* current generation count */
        le32 offset;        /* configuration offset */
        le32 length;        /* number of bytes in affected range */
        u8   data[];        /* optional configuration data */
};
\end{lstlisting}

The conditions that require \msgref{EVENT_CONFIG}, and its generation and
\field{offset}/\field{length} semantics, are defined by
\ref{sec:Virtio Transport Options / Virtio Over Messages / Basic Concepts / Configuration Semantics Profiles}
and
\ref{sec:Virtio Transport Options / Virtio Over Messages / Device Operation / Device Notifications}.

\msgdef{EVENT_AVAIL}

\msgref{EVENT_AVAIL} notifies the device that a virtqueue has new buffers. The
message includes the virtqueue index in \field{vq_index} and, if
VIRTIO\_F\_NOTIFICATION\_DATA is negotiated, a combined next-offset/next-wrap
field in \field{next_offset} so the device can jump directly to the updated
position. If VIRTIO\_F\_NOTIFICATION\_DATA is not negotiated, \field{next_offset}
is reserved and is set to zero. This transport does not support
VIRTIO\_F\_NOTIF\_CONFIG\_DATA; \field{vq_index} always identifies the target
virtqueue by index. When VIRTIO\_F\_NOTIFICATION\_DATA is negotiated,
\field{next_offset} encodes core virtio notification data:
\field{next_offset}[30:0] carries core \field{next_off}, and
\field{next_offset}[31] carries core \field{next_wrap}. For split virtqueues,
\field{next_off} is derived from split-ring available-index progression and
\field{next_wrap} is zero. For packed virtqueues, \field{next_off} and
\field{next_wrap} are derived from packed-ring notification-data values as
defined by core virtio.

\begin{lstlisting}
struct virtio_msg_event_avail {
        le32 vq_index;    /* virtqueue index */
        le32 next_offset; /* bits[30:0] offset, bit[31] wrap; reserved if unused */
};
\end{lstlisting}

\drivernormative{\paragraph}{EVENT\_AVAIL}{Virtio Transport Options / Virtio Over Messages / Transport Messages / VIRTIO_MSG_EVENT_AVAIL / Driver}
\begin{itemize}
  \item A driver MUST identify the target virtqueue in \field{vq_index}.
  \item If VIRTIO\_F\_NOTIFICATION\_DATA has not been negotiated, a driver MUST
        set \field{next_offset} to zero in \msgref{EVENT_AVAIL} messages. If it
        has been negotiated, the driver MUST encode core \field{next_off} in
        \field{next_offset}[30:0] and core \field{next_wrap} in
        \field{next_offset}[31], using split-queue or packed-queue derivation
        as defined by core virtio.
\end{itemize}

\busnormative{\paragraph}{EVENT\_AVAIL}{Virtio Transport Options / Virtio Over Messages / Transport Messages / VIRTIO_MSG_EVENT_AVAIL / Bus}
\begin{itemize}
  \item If a bus implementation cannot preserve
        VIRTIO\_F\_NOTIFICATION\_DATA semantics for \msgref{EVENT_AVAIL}, it
        MUST prevent negotiation of VIRTIO\_F\_NOTIFICATION\_DATA.
\end{itemize}

\msgdef{EVENT_USED}

\msgref{EVENT_USED} notifies the driver that buffers on a particular virtqueue
have been consumed. The payload carries only the queue index; the driver uses
its used ring to determine which descriptors completed. If polling-only mode is
selected by the endpoints, \msgref{EVENT_USED} is not a semantic requirement
for completion observation. Polling-only selection and the requirement to
originate \msgref{EVENT_USED} when event-driven completion is used are defined
in
\ref{sec:Virtio Transport Options / Virtio Over Messages / Device Operation / Device Notifications}.

\begin{lstlisting}
struct virtio_msg_event_used {
        le32 vq_index; /* virtqueue index */
};
\end{lstlisting}

\devicenormative{\paragraph}{EVENT\_USED}{Virtio Transport Options / Virtio Over Messages / Transport Messages / VIRTIO_MSG_EVENT_USED / Device}
\begin{itemize}
  \item A device MUST identify the virtqueue whose buffers have been consumed
        in \field{vq_index}.
\end{itemize}


\subsection{Bus Messages}\label{sec:Virtio Transport Options / Virtio Over Messages / Bus Messages}

The standard bus messages defined here provide optional mechanisms for device
discovery, hotplug notifications, and bus-level liveness checks. A bus binding
may use any subset of these messages, including none, and may define alternative
mechanisms for these functions. Support for a standard bus message does not
require support for dynamic device addition or removal.

\subsubsection{Overview}
\label{sec:Virtio Transport Options / Virtio Over Messages / Bus Messages / Overview}

Each bus implementation chooses which of these messages to support. Some may
implement all of them for a fully message-based workflow, while others may rely
on out-of-band enumeration and use only the subset that fits their needs.

\paragraph{Message IDs and initiating side}

In the table below, the initiating side identifies which side of the bus
message exchange starts the message, not the local software component that
serializes it on a particular implementation.

\begin{tabular}{|l|l|l|}
\hline
Name & ID & Initiating side \\
\hline
\hline
Reserved                    & 0x0  &        \\
\hline
Reserved                    & 0x1  &        \\
\hline
\busref{GET_DEVICES}        & 0x2  & Driver side \\
\hline
\busref{PING}               & 0x3  & Either side \\
\hline
\busref{EVENT_DEVICE}       & 0x40 & Device side \\
\hline
\end{tabular}

Bus-message ID bit encoding is defined in
\ref{sec:Virtio Transport Options / Virtio Over Messages / Basic Concepts / Common Message Format}.

\paragraph{Bus Specific Messages}
\label{sec:Virtio Transport Options / Virtio Over Messages / Bus Messages / Bus Specific Messages}

A range of message IDs are reserved for use by the specific bus
implementation. These messages can be used for any implementation-specific
usage. Example usage could include:

\begin{itemize}
  \item Configuration of out-of-band notification methods
  \item Setup shared memory regions to be used for buffers or virtqueues
  \item Declare bus-specific error conditions
  \item Provide extra debug or logging information
\end{itemize}

\busdef{GET_DEVICES}

A bus implementation uses this request to enumerate device numbers exposed by a
peer bus instance. The response describes a window of device numbers and
signals which entries are present.

\begin{lstlisting}
struct virtio_bus_msg_get_devices_req {
        le16 offset; /* starting device number */
        le16 count;  /* number of device-number slots requested */
};

struct virtio_bus_msg_get_devices_resp {
        le16 offset;      /* echoed starting device number */
        le16 next_offset; /* 0 or suggested start for next query */
        le16 count;       /* number of slots returned */
        u8   bitmap[];    /* ((count + 7) / 8) bytes, LSB-first packing */
};
\end{lstlisting}

The (\field{offset}, \field{count}) tuple defines a window of \field{count}
consecutive device numbers beginning at \field{offset}. A request with
\field{count}=0 is valid and requests an empty window. The number of present
devices equals the number of set bits in the first \field{count} positions of
\field{bitmap}. The \field{bitmap} size in bytes is
\textbf{(\field{count} + 7) / 8}. The response \field{count} MUST NOT exceed
the request \field{count}. Responders SHOULD return the requested
\field{count} unless constrained
(e.g., by maximum message size) or when trailing parts of the requested window
contain no present device numbers; in those cases they MAY return a smaller
value, and the \field{bitmap} covers the reduced window. A response with
\field{count}=0 returns an empty \field{bitmap}; this is valid for a request
with \field{count}=0 and also when no devices are present in the returned
window.

Example: a request with \field{offset}=0, \field{count}=16 might produce
\field{bitmap}=0b00100101 00000000 and \field{next_offset}=16. That indicates
devices 0, 2, and 5 are present within the 0-15 window and suggests continuing
at device number 16.

Drivers can use this message to enumerate device numbers. Each window can be
treated as a snapshot, advanced using \field{next_offset}, and validated with
\msgref{GET_DEVICE_INFO} before other transport messages are issued. A response
with \field{next_offset}=0 is terminal. Otherwise, \field{next_offset}
recommends a later starting device number and advances enumeration beyond the
request \field{offset}.

\busnormative{\paragraph}{GET\_DEVICES}{Virtio Transport Options / Virtio Over Messages / Bus Messages / GET_DEVICES / Bus}
\begin{itemize}
  \item The responding bus implementation MUST echo the requested
        \field{offset} in its response.
  \item The responding bus implementation MUST interpret \field{count} as a
        number of device-number slots and MUST return a \field{bitmap} of
        \textbf{(\field{count} + 7) / 8} bytes, packed least-significant-bit
        first.
  \item A request with \field{count}=0 is valid. In response to such a request,
        the responding bus implementation MUST return \field{count}=0 and an
        empty \field{bitmap}.
  \item The responding bus implementation MUST NOT return a \field{count} value
        greater than the \field{count} value from the request.
  \item If \field{count} is not a multiple of 8, the responder MUST set unused
        bits in the last \field{bitmap} byte (positions at or above
        \field{count}) to zero.
  \item For a terminal response, the responding bus implementation MUST set
        \field{next\_offset} to 0.
  \item For a nonterminal response, the responding bus implementation MUST set
        \field{next\_offset} to a value strictly greater than the request
        \field{offset}.
  \item Responders SHOULD return the requested \field{count} unless constrained
        (e.g., by maximum message size) or when trailing parts of the requested
        window contain no present device numbers; if a smaller \field{count} is
        returned, the \field{bitmap} MUST still represent the window defined by
        the echoed \field{offset} and \field{count}.
  \item If no devices are present in the requested window, a responder MAY
        return \field{count}=0 with an empty \field{bitmap}.
\end{itemize}

\busdef{PING}

\busref{PING} checks whether the peer bus endpoint responds. It can be sent by
either side of the bus; the response echoes the 32-bit data value from the
request. This check can help detect failures when independently running systems
communicate and one stops without notifying its peer, including during periods
without other message traffic.

\begin{lstlisting}
struct virtio_bus_msg_ping_req {
        le32 data; /* arbitrary value chosen by the sender */
};

struct virtio_bus_msg_ping_resp {
        le32 data; /* echoed request data */
};
\end{lstlisting}

Support for \busref{PING} is optional. Its use, sending interval, timeout, and
recovery policy are determined by the bus binding and deployment, normally
within the bus implementation. A successful response establishes peer bus
endpoint responsiveness; it does not establish the health or presence of
individual devices.

\busnormative{\paragraph}{PING}{Virtio Transport Options / Virtio Over Messages / Bus Messages / PING / Bus}
\begin{itemize}
  \item The responding bus implementation MUST echo the 32-bit data field from
        the request exactly in the \busref{PING} response.
\end{itemize}

\busdef{EVENT_DEVICE}

\busref{EVENT_DEVICE} signals device-level hotplug changes or changes requiring
bus-specific handling as a bus-originated control-plane event. The bus instance
on the driver side receives these indications from the peer bus instance; the
table below lists the defined device bus states.

\begin{lstlisting}
struct virtio_bus_msg_event_device {
        le16 device_number;
        le16 device_bus_state; /* see table below */
};
\end{lstlisting}

\begin{tabular}{|l|l|p{7cm}|}
\hline
Constant & Value & Meaning \\
\hline
DEVICE\_BUS\_STATE\_NO\_DATA  & 0x0000 & Change requiring bus-specific handling. \\
DEVICE\_BUS\_STATE\_ADDED    & 0x0001 & Device is present and ready. \\
DEVICE\_BUS\_STATE\_REMOVED  & 0x0002 & Device is no longer present. \\
Reserved                    & 0x0003-0x7FFF & Reserved for standard use. \\
Implementation-defined      & 0x8000-0xFFFF & MAY be used by the bus implementation. \\
\hline
\end{tabular}

DEVICE\_BUS\_STATE\_NO\_DATA indicates that a change has occurred without
specifying its nature or identifying individual devices. The bus binding
defines the required handling. For example, a bus binding may require the
driver-side bus to reset the bus connection, remove all devices from its local
view, and re-enumerate the device set. The \field{device\_number} field is zero
and does not identify a device in this case; the event retains the same payload
structure. The removal-safety and device-number assignment requirements
continue to apply.

Upon receiving DEVICE\_BUS\_STATE\_ADDED, drivers typically probe the device
with \msgref{GET_DEVICE_INFO} before continuing initialization.
DEVICE\_BUS\_STATE\_REMOVED reports that the device is no longer accessible
for transport-message processing. Removal handling follows
\ref{sec:Virtio Transport Options / Virtio Over Messages / Device Operation / Hotplug and Removal}.
Transport requests that cannot be routed after removal are subject to
\ref{busnormative:Virtio Transport Options / Virtio Over Messages / Basic Concepts / Relationship between bus and transport / Transport Message Forwarding}.
One-way events have no request completion outcome and follow the event delivery
rules in
\ref{sec:Virtio Transport Options / Virtio Over Messages / Basic Concepts / Error Signaling}.

\busnormative{\paragraph}{EVENT\_DEVICE}{Virtio Transport Options / Virtio Over Messages / Bus Messages / EVENT_DEVICE / Bus}
\begin{itemize}
  \item A bus implementation that sends \busref{EVENT_DEVICE} MUST set
        \field{device\_number} to a valid bus device number exposed to the
        driver side, except for DEVICE\_BUS\_STATE\_NO\_DATA.
  \item A bus implementation that sends \busref{EVENT_DEVICE} with
        DEVICE\_BUS\_STATE\_NO\_DATA MUST set \field{device\_number} to zero.
  \item A bus implementation that uses DEVICE\_BUS\_STATE\_NO\_DATA MUST
        define the required handling in its bus binding.
  \item A bus implementation that receives \busref{EVENT_DEVICE} with
        DEVICE\_BUS\_STATE\_NO\_DATA MUST ignore \field{device\_number} and
        perform the handling defined by its bus binding.
  \item A bus implementation that sends \busref{EVENT_DEVICE} with
        DEVICE\_BUS\_STATE\_ADDED MUST do so only when the referenced device is
        accessible for transport-message processing.
  \item A bus implementation that sends \busref{EVENT_DEVICE} with
        DEVICE\_BUS\_STATE\_REMOVED MUST do so when the referenced device is no
        longer available for transport-message processing.
\end{itemize}
